\section{Quadratic forms with hub coordinates}\label{sec:hubs}

Let $X_0,X_1,\ldots,X_m$ be independent real random variables with
probability densities bounded by one, and write $X'=(X_1,\ldots,X_m)$.
For orthonormal vectors $w,v\in\mathbb R^m$ put
$L=\langle w,X'\rangle$ and $V=\langle v,X'\rangle$.
The basic family in this section is
\begin{equation}\label{eq:hub-normal}
 Q=aX_0^2+2X_0(L+\beta)+\lambda V^2
       +2L\langle q,X'\rangle+\sum_{i=1}^m\phi_i(X_i),
 \qquad \phi_i(x)=d_ix^2+b_ix+c_i,
\end{equation}
where $\lambda\ne0$.  Unless a stronger condition is stated, assume
\begin{equation}\label{eq:hub-frame}
 w_i^2+v_i^2\le\tfrac12\qquad(1\le i\le m).
\end{equation}
The coefficients of the one-coordinate polynomials are unrestricted.
We use the Fourier convention
$\widehat\mu(t)=\int e^{itx}\,\mu(dx)$, and all logarithms are natural.

The one-coordinate term $X_0$ is a hub: its coupling to the remaining
coordinates is carried by the single linear statistic $L$.  The estimates
below distinguish the independent hub law from the dependence between
$L$, $V$, and the other linear statistics of $X'$.

\subsection{Projection and Fourier estimates}

We use the following consequence of the density bound for projections of
product measures \citep{LPP2016,RV2015}: if $T$ is an $r\times m$
matrix of rank $r$, then
\begin{equation}\label{eq:hub-projection}
 \|p_{TX'}\|_\infty\le
       \frac{2^{r/2}}{\sqrt{\det(TT^T)}}.
\end{equation}
The bound for an orthonormal projection is $2^{r/2}$.  No independence
between the projected coordinates is asserted.

\begin{lemma}[Modulated frame integral]\label{lem:hub-frame-integral}
Let $f\le1$ be a probability density and let $g$ be a complex function
with $|g|\le f$.  For $p\ge2$,
\begin{equation}\label{eq:hub-scalar-lp}
 \int_{\mathbb R}|\widehat g(t)|^p\,dt\le\frac{4\pi}{\sqrt p}.
\end{equation}
Let $e_1,\ldots,e_r$ be orthonormal vectors in $\mathbb R^m$, set
$u_i=(e_{1i},\ldots,e_{ri})$, and assume $|u_i|^2\le1/2$.
For complex functions $g_i$ dominated by the coordinate densities and
arbitrary real shifts $s_i$,
\begin{equation}\label{eq:hub-frame-int}
 \int_{\mathbb R^r}\prod_i
   |\widehat g_i(\langle u_i,z\rangle+s_i)|\,dz\le(4\pi)^r.
\end{equation}
\end{lemma}

\begin{proof}
For an integer $k\ge1$, Plancherel and convolution give
\[
 \int|\widehat g|^{2k}
 =2\pi(g*\widetilde g)^{*k}(0)
 \le2\pi(f*\widetilde f)^{*k}(0),
 \qquad \widetilde g(x)=\overline{g(-x)}.
\]
The density on the right is that of a sum of $2k$ independent variables
whose densities are bounded by one.  Applying
\eqref{eq:hub-projection} to this one-dimensional linear map gives a
bound $1/\sqrt k$ for that density.  Since $|\widehat g|\le1$, choose
$k=\lfloor p/2\rfloor$.  For $p\ge4$, $k\ge p/4$, and for
$2\le p<4$ Plancherel gives the required bound directly.

Set $c_i=|u_i|^2$ and omit the zero rows.  The vectors
$u_i/\sqrt{c_i}$ with weights $c_i$ form geometric
Brascamp--Lieb data, since $\sum_i u_i u_i^T=I_r$.
The geometric Brascamp--Lieb inequality
\citep{BL1976,Ball1986,Barthe1998}, including arbitrary translations of
the one-dimensional arguments, bounds the integral in
\eqref{eq:hub-frame-int} by
\[
 \prod_i\left(c_i^{-1/2}
          \int|\widehat g_i(t)|^{1/c_i}\,dt\right)^{c_i}
 \le(4\pi)^{\sum_i c_i}=(4\pi)^r.
\]
Here $1/c_i\ge2$, so \eqref{eq:hub-scalar-lp} applies.
\end{proof}

For $\ell\in\mathbb R$ define the frozen polynomial
\[
 Y_\ell=\lambda V^2+2\ell\langle q,X'\rangle
                         +\sum_i\phi_i(X_i).
\]
For each fixed $\ell_0$, let $\mu_{\ell_0,\ell}$ denote the
unnormalized slice of the joint law of $(L,Y_{\ell_0})$ at $L=\ell$.
Its total mass is $\rho_L(\ell)$, not one.

\begin{lemma}[Frozen slices]\label{lem:hub-frozen}
For every fixed $\ell_0$, for almost every $\ell$, and every $t\ne0$,
\begin{equation}\label{eq:hub-slice-fourier}
 |\widehat\mu_{\ell_0,\ell}(t)|
 \le\min\left\{\rho_L(\ell),
                   \frac{4\sqrt\pi}{\sqrt{|t\lambda|}}\right\}.
\end{equation}
In particular, for every interval $I$,
\begin{equation}\label{eq:hub-slice-interval}
 \mu_{\ell_0,\ell}(I)
       \le2\pi^2\sqrt{\frac{|I|}{|\lambda|}}.
\end{equation}
\end{lemma}

\begin{proof}
Write $Y_{\ell_0}=\lambda V^2+\sum_i(d_iX_i^2+c_i'X_i)+c'$.
For fixed $t\ne0$ the Fourier transform in $L$ of the complex measure
$\mathbb E[e^{itY_{\ell_0}};L\in\cdot]$ is
\[
 \Psi(z,t)=\mathbb E e^{izL+itY_{\ell_0}}.
\]
The Fresnel identity, first with Gaussian damping, gives
\[
 |\Psi(z,t)|\le(4\pi|t\lambda|)^{-1/2}
  \int\prod_i|\widehat g_i(zw_i+uv_i+tc_i')|\,du,
 \qquad g_i(x)=f_i(x)e^{itd_ix^2}.
\]
Lemma~\ref{lem:hub-frame-integral} bounds the integral in $(z,u)$ by
$(4\pi)^2$.  Fourier inversion in $z$ therefore gives the second term
in \eqref{eq:hub-slice-fourier}.  The first follows from the total mass
of the slice.  Disintegration first identifies this transform for almost
every pair $(\ell,t)$; continuity in $t$ identifies it for every $t$ on
a common full-measure set of $\ell$.

If $I$ has center $y_0$ and length $2b$, the function
\[
 K(y)=\left(\frac{\sin(\pi(y-y_0)/(4b))}
                         {\pi(y-y_0)/(4b)}\right)^2
\]
satisfies $K\ge(8/\pi^2)\mathbf1_I$; its Fourier transform is supported
on $|t|\le\pi/(2b)$ and has modulus at most $4b$.
Parseval and \eqref{eq:hub-slice-fourier} give
\[
 \mu(I)\le\frac{\pi^2}{8}\frac{4b}{2\pi}
   \int_{|t|\le\pi/(2b)}\frac{4\sqrt\pi}{\sqrt{|t\lambda|}}\,dt
 =2\pi^2\sqrt{\frac{2b}{|\lambda|}}.
\]
All the Fourier manipulations follow first with a positive Gaussian
damping; the displayed integrable bounds justify its removal.
\end{proof}

\begin{remark}[Freezing and exceptional slices]\label{rem:hub-freezing}
An estimate valid for every fixed $\ell_0$ and almost every $\ell$ cannot
be substituted on the diagonal $\ell_0=\ell$ without justification.
We use the following discretization whenever the slice parameter varies.
Partition the real line into cells $C_j$ of width $h$, choose
$\ell_j\in C_j$, and replace $Y_L$ by $Y_{\ell_j}$ on
$\{L\in C_j\}$.  Then
$Y_{\ell_j}-Y_L=2(\ell_j-L)\langle q,X'\rangle\to0$ almost surely.
For a bounded kernel $G(\ell,y)$ continuous off the null set
$\ell=-\beta$, dominated convergence gives
\begin{equation}\label{eq:hub-freezing-limit}
 \mathbb E G(L,Y_L)
 =\lim_{h\downarrow0}\sum_j\int_{C_j}\int
       G(\ell_j,y)\,\mu_{\ell_j,\ell}(dy)\,d\ell.
\end{equation}
For each partition only countably many frozen parameters occur, so the
slice bounds apply simultaneously.  The analogous construction with
rectangular cells applies to several hub partners.  Riemann sums used
below can first be restricted to bounded sets; tails are controlled by
the total mass.  This also justifies the formulas for arbitrary
unbounded coordinate laws.
\end{remark}

\subsection{Real-variable estimates for one hub}

\begin{lemma}[A profile criterion]\label{lem:hub-profile}
Assume $a=0$.  Suppose that for every fixed $\ell_0$, almost every
$\ell$, and every Borel set $S$,
\[
 \mu_{\ell_0,\ell}(S)\le K|S|^\theta,
 \qquad 0<\theta\le1.
\]
Then
\begin{equation}\label{eq:hub-profile-bound}
 \|p_Q\|_\infty\le\frac12+\frac{2^\theta K}{\theta}.
\end{equation}
More generally, a concave nondecreasing profile $\omega$ with
$\omega(0)=0$ gives
$\|p_Q\|_\infty\le\frac12+\int_0^2\omega(u)\,du/u$ whenever
that integral is finite.
\end{lemma}

\begin{proof}
Fix a level $s$ and a window radius $\delta>0$.  With $t=\ell+\beta$ put
\[
 G(\ell,y)=\int f_0(x)
            \mathbf1_{\{|2tx+y-s|\le\delta\}}\,dx.
\]
For $t\ne0$, $0\le G\le\delta/|t|$, and $\int G(\ell,y)\,dy=2\delta$.
If $M=\delta/|t|$, layer cake and Jensen's inequality give
\[
 \int G\,d\mu
  \le K\int_0^M |\{G>u\}|^\theta\,du
  \le KM(2\delta/M)^\theta
  =K\delta\,2^\theta|t|^{\theta-1}.
\]
For $|t|\ge1$ use $G\le\delta$ and the slice mass
$\rho_L(\ell)$, whose integral is one.  Apply
\eqref{eq:hub-freezing-limit}, integrate the displayed bound over
$|t|<1$, divide by $2\delta$, and let $\delta\downarrow0$.
This proves \eqref{eq:hub-profile-bound}.  The same layer-cake argument
gives $M\omega(2\delta/M)$ for a concave profile and yields the final
assertion.  The resulting uniform interval bounds imply absolute
continuity and the stated essential density bound.
\end{proof}

\begin{theorem}[An interval-supported hub]\label{thm:hub-interval}
Let $a=0$ and suppose $X_0$ is supported in an interval of length $D$.
Then every form \eqref{eq:hub-normal} satisfies
\begin{equation}\label{eq:hub-interval-bound}
 \|p_Q\|_\infty\le
 \frac32(4\sqrt2\pi^2)^{2/3}(D/|\lambda|)^{1/3}
 <22(D/|\lambda|)^{1/3}.
\end{equation}
\end{theorem}

\begin{proof}
Let $J$ be the supporting interval and replace $f_0$ by $\mathbf1_J$.
For $t=\ell+\beta\ne0$ the resulting window kernel is
\[
 G(\ell,y)=\frac{|(s-y+[-\delta,\delta])\cap2tJ|}{2|t|}.
\]
Three estimates follow by Fubini and
Lemma~\ref{lem:hub-frozen}:
\begin{align*}
 \int G\,d\mu&\le (\delta/|t|)\rho_L(\ell),\\
 \int G\,d\mu&\le (\delta/|t|)2\pi^2
                  \sqrt{(2|t|D+2\delta)/|\lambda|},\\
 \int G\,d\mu&\le 2\pi^2D\sqrt{2\delta/|\lambda|}.
\end{align*}
Use the first on $|t|\ge T$ and the minimum of the other two on
$|t|<T$, with the freezing procedure of
Remark~\ref{rem:hub-freezing}.  After division by $2\delta$, the range
$|t|<\delta$ tends to zero.  On $\delta\le|t|<T$, dominated
convergence applies with an integrable multiple of $|t|^{-1/2}$.
It follows that, uniformly in $s$,
\[
 \limsup_{\delta\downarrow0}
 \frac{\mathbb P(|Q-s|\le\delta)}{2\delta}
 \le\frac1{2T}+4\sqrt2\pi^2\sqrt{DT/|\lambda|}.
\]
Choose $T=(4\sqrt2\pi^2)^{-2/3}(D/|\lambda|)^{-1/3}$.
\end{proof}

The next result needs no spread condition.  Its only use of the product
law is the two-dimensional projection bound.

\begin{theorem}[A transverse square with an independent hub]
\label{thm:hub-pure}
Let $w,v$ be orthonormal and let $v_c,F_c$ be arbitrary measurable
real functions.  Set
\[
 Q=2X_0(L+\beta)+\lambda(V-v_c(L))^2+F_c(L),
 \qquad \mu=|\lambda|>0.
\]
Then
\begin{equation}\label{eq:hub-pure-bound}
 \|p_Q\|_\infty\le
 \begin{cases}
 2\sqrt2+\sqrt2\log(1/3+4\sqrt2/\mu),&0<\mu\le6\sqrt2,\\
 (192/\mu)^{1/3},&\mu\ge6\sqrt2.
 \end{cases}
\end{equation}
\end{theorem}

\begin{proof}
For fixed $s$, integrating out $X_0$ gives
\[
 p_Q(s)=\int\rho_{L,V}(\ell,v)
 \frac{f_0((s-\lambda(v-v_c(\ell))^2-F_c(\ell))/(2t))}
      {2|t|}\,d\ell\,dv,
 \qquad t=\ell+\beta.
\]
Denote the fraction by $K(\ell,v)$.  It is at most $1/(2|t|)$.
Changing variables on the two branches of the square and using
$f_0\le1$, $\int f_0=1$, gives
\[
 \int K(\ell,v)\,dv\le\sqrt2(\mu|t|)^{-1/2}.
\]
Indeed the remaining integral is a one-sided negative half moment of
$f_0$, bounded by $\int_0^1u^{-1/2}\,du=2$.
Since $\rho_{L,V}\le2$ and $\rho_L\le\sqrt2$, the contribution at $t$
is bounded by $g(t)/(2|t|)$ for some nonnegative function satisfying
\[
 \int g\le1,\qquad
 g(t)\le M(t):=\min\{\sqrt2,4\sqrt2\mu^{-1/2}\sqrt{|t|}\}.
\]
The weight $1/(2|t|)$ decreases with $|t|$.  Thus its maximal integral
under these constraints is obtained by taking $g=M$ on
$[-\tau,\tau]$, where $\int_{-\tau}^{\tau}M=1$.
The envelope switches at $t_* =\mu/16$, and its mass on
$[-t_*,t_*]$ is $\sqrt2\mu/12$.
If this mass is at most one, the inner contribution is $2\sqrt2$ and
the outer contribution is $\sqrt2\log(\tau/t_*)$, with
$\tau/t_*=1/3+4\sqrt2/\mu$.
Otherwise
$\tau=(3\sqrt\mu/(16\sqrt2))^{2/3}$, and the integral is
$8\sqrt2\mu^{-1/2}\sqrt\tau=(192/\mu)^{1/3}$.
The argument can equivalently be applied to interval probabilities,
so no prior existence of $p_Q$ is required.
\end{proof}

\subsection{A nonzero transverse cross term}

Write
\[
 q=q_ww+q_vv+\varepsilon e,\qquad
 e\perp w,v,\quad |e|=1,\quad\varepsilon>0,
\]
and set $P=\langle e,X'\rangle$.

\begin{theorem}[Nonzero leakage]\label{thm:hub-leak}
Assume $w_i^2+v_i^2+e_i^2\le1/2$ for every light coordinate.  For
\eqref{eq:hub-normal}, including an arbitrary hub diagonal $a$,
\begin{align}
 |\mathbb E e^{itQ}|&\le
 4\sqrt2\pi^{3/2}(\varepsilon|\lambda|)^{-1/2}|t|^{-3/2},
 \label{eq:hub-leak-fourier}\\
 \|p_Q\|_\infty&\le
 3\,2^{5/3}(\varepsilon|\lambda|)^{-1/3}.
 \label{eq:hub-leak-density}
\end{align}
\end{theorem}

\begin{proof}
Complete the square in $V$:
$V'=V+(q_v/\lambda)L$ and
$\sigma=2q_w-q_v^2/\lambda$.
After moving $-2\varepsilon\beta P$ into the linear terms, write
\[
 Q=aX_0^2+2X_0(L+\beta)+\sigma L^2+
   2\varepsilon(L+\beta)P+W,
 \quad
 W=\lambda V'^2-2\varepsilon\beta P+\sum_i\phi_i(X_i).
\]
Let $g(\omega)=\int f_0(x)e^{ita x^2+i\omega x}\,dx$ and
\[
 H(\ell,p)=g(2t(\ell+\beta))
                   e^{it\sigma\ell^2+2it\varepsilon(\ell+\beta)p}.
\]
Let $n(\ell,p)$ be the density of
$\mathbb E[e^{itW};(L,P)\in\cdot]$.  Its Fourier transform is
integrable by the Fresnel identity and the three-dimensional frame
integral.  Fourier transformation of $H$ in $p$ pins
$\ell+\beta=-\gamma/(2t\varepsilon)$; consequently
\[
 |\widehat H(z,\gamma)|=
      \frac{\pi}{|t|\varepsilon}|g(-\gamma/\varepsilon)|.
\]
This identity is interpreted with Gaussian damping, or as a tempered
distribution followed by the indicated integrable pairing.
Put $G(\gamma)=\int|\widehat n(z,\gamma)|\,dz$.
Fresnel linearization in $V'$ and translation of the $z$ variable give
\[
 \|G\|_\infty\le(4\pi|t\lambda|)^{-1/2}(4\pi)^2,
 \qquad
 \|G\|_1\le(4\pi|t\lambda|)^{-1/2}(4\pi)^3.
\]
Thus $\|G\|_2\le(4\pi|t\lambda|)^{-1/2}(4\pi)^{5/2}$.
Plancherel gives
$\int|g(\gamma/\varepsilon)|^2\,d\gamma\le2\pi\varepsilon$.
Parseval and Cauchy--Schwarz yield
\[
 |\mathbb E e^{itQ}|
 \le\frac{1}{4\pi|t|\varepsilon}
       (2\pi\varepsilon)^{1/2}
       (4\pi|t\lambda|)^{-1/2}(4\pi)^{5/2},
\]
which is \eqref{eq:hub-leak-fourier}.
Finally, $\int_{\mathbb R}\min(1,B|t|^{-3/2})\,dt=6B^{2/3}$.
Fourier inversion gives \eqref{eq:hub-leak-density}.
\end{proof}

\begin{example}[A frozen-hub obstruction]\label{ex:hub-frozen-obstruction}
The nonzero-leakage estimate cannot be made uniform by bounding every
conditional density obtained on fixing $X_0=x$.  Fix $x=r>0$, let the
light vector have law $X'=-(r/\varepsilon)e+G$, where the coordinates
of $G$ are independent standard Gaussians, and take no other terms.
Then the conditional quadratic form is
\[
 \lambda\langle v,G\rangle^2+
       2\varepsilon\langle w,G\rangle\langle e,G\rangle.
\]
For $\lambda>0$, let $Z_1,Z_2,Z_3$ be its three independent standard
Gaussian coordinates.  The density of $2Z_1Z_2$ is positive and
continuous on any compact interval avoiding zero.  In the convolution
formula for $\lambda Z_3^2+2\varepsilon Z_1Z_2$ at zero, restrict
$\lambda Z_3^2$ to $[\varepsilon,2\varepsilon]$.
The square density is bounded below by a constant times
$(\lambda u)^{-1/2}$ there when $\varepsilon\le\lambda$; the second
density is bounded below by a constant times $\varepsilon^{-1}$.
Integration gives a lower bound
$c(\lambda\varepsilon)^{-1/2}$.
All input densities are bounded by one.  This is an obstruction to
uniform frozen-hub estimates, not a counterexample to the full quadratic
form after averaging the independent hub.
\end{example}

\subsection{Near-linear small-ball bounds and symmetric-window recovery}

\begin{theorem}[Double-logarithmic small balls]\label{thm:hub-loglog-smallball}
Assume $a=0$ and \eqref{eq:hub-frame}.  No restriction is imposed on
$q$ beyond \eqref{eq:hub-normal}.  For every $s\in\mathbb R$ and
$\delta>0$,
\begin{equation}\label{eq:hub-loglog-smallball}
 \frac{\mathbb P(|Q-s|\le\delta)}{2\delta}
 \le A_\delta:=\frac12+2\sqrt2+
                 \sqrt2\log_+(2B_\delta^2),
\end{equation}
where
\begin{equation}\label{eq:hub-bdelta}
 B_\delta=\frac{4+\sqrt{4/\pi+32\log_+(|\lambda|/\delta)}}
                     {2\sqrt{|\lambda|}}.
\end{equation}
\end{theorem}

\begin{proof}
Use the window kernel in the proof of Lemma~\ref{lem:hub-profile}.
For fixed $t=\ell+\beta\ne0$, it is
$G(\ell,\cdot)=\mathbf1_{[-\delta,\delta]}*g$, where
$g(u)=f_0((s-u)/(2t))/(2|t|)$.  Hence
\[
 \widehat G(\eta)=\frac{2\sin(\eta\delta)}{\eta}\widehat g(\eta),
 \qquad \int|\widehat g|^2\le\frac\pi{|t|}.
\]
For every applicable frozen slice, Lemma~\ref{lem:hub-frozen} gives
\[
 \int_{|\eta|\le1/\delta}|\widehat\mu(\eta)|^2\,d\eta
 \le\frac{4+32\pi\log_+(|\lambda|/\delta)}{|\lambda|}.
\]
To see this, use $|\widehat\mu|^2\le2$ below $1/|\lambda|$
and $|\widehat\mu|^2\le16\pi/(|\eta\lambda|)$ above it.
Parseval and Cauchy--Schwarz on the regions below and above
$1/\delta$ now give
\begin{align*}
 \int G\,d\mu
 &\le\frac1{2\pi}\left(
 2\delta\int_{|\eta|\le1/\delta}|\widehat g||\widehat\mu|
 +\int_{|\eta|>1/\delta}\frac2{|\eta|}
                          |\widehat g||\widehat\mu|\right)\\
 &\le\delta(|\lambda t|)^{-1/2}
          \left(\sqrt{4/\pi+32\log_+(|\lambda|/\delta)}+4\right).
\end{align*}
The high-frequency calculation uses
$\int_{|\eta|>1/\delta}|\eta|^{-3}\,d\eta=\delta^2$.
The alternative bound is
$\int G\,d\mu\le\sqrt2\delta/|t|$.
For $|t|\ge1$ retain the slice mass and obtain a contribution at most
$\delta$.  Formula~\eqref{eq:hub-freezing-limit} therefore implies
\[
 \frac{\mathbb P(|Q-s|\le\delta)}{2\delta}
 \le\frac12+\int_{|t|<1}
       \min\left\{\frac{\sqrt2}{2|t|},
                         B_\delta|t|^{-1/2}\right\}\,dt.
\]
If $2B_\delta^2\ge1$, split at
$|t|=1/(2B_\delta^2)$; the integral is
$2\sqrt2+\sqrt2\log(2B_\delta^2)$.
Otherwise it equals $4B_\delta\le2\sqrt2$.
\end{proof}

\begin{lemma}[Recovery from a symmetric window]\label{lem:hub-window-recovery}
Suppose a real random variable $Z$ satisfies
$|\mathbb E e^{itZ}|\le B|t|^{-3/2}$ for $t\ne0$.
Then it has a continuous bounded density and, for every $h>0$,
\begin{equation}\label{eq:hub-window-recovery}
 \|p_Z\|_\infty\le
 \sup_s\frac{\mathbb P(|Z-s|\le h)}{2h}
       +\frac{2\sqrt{2\pi}}{3\pi}B\sqrt h.
\end{equation}
\end{lemma}

\begin{proof}
The integrable characteristic function gives the density by Fourier
inversion.  Let $U_h=\mathbf1_{[-h,h]}/(2h)$.
Its Fourier transform is $\sin(ht)/(ht)$, which is at most one.
Consequently
\[
 |p_Z(s)-(p_Z*U_h)(s)|
 \le\frac{B\sqrt h}{\pi}
   \int_0^\infty u^{-3/2}\left(1-\frac{\sin u}{u}\right)\,du.
\]
Since $\sin u/u=\int_0^1\cos(ru)\,dr$, Tonelli and an integration by
parts give
\begin{align*}
 \int_0^\infty u^{-3/2}\left(1-\frac{\sin u}{u}\right)\,du
 &=\int_0^1\int_0^\infty u^{-3/2}(1-\cos(ru))\,du\,dr\\
 &=\sqrt{2\pi}\int_0^1\sqrt r\,dr
 =\frac23\sqrt{2\pi}.
\end{align*}
For completeness, the inner integral at $r=1$ is
$2\int_0^\infty u^{-1/2}\sin u\,du=\sqrt{2\pi}$;
this identity follows by exponential damping and the Gamma integral.
Finally $(p_Z*U_h)(s)=\mathbb P(|Z-s|\le h)/(2h)$.
\end{proof}

\begin{corollary}[Improved double-logarithmic density estimate]
\label{cor:hub-loglog-density}
Under the assumptions of Theorem~\ref{thm:hub-leak}, with $a=0$,
\begin{equation}\label{eq:hub-loglog-recovery}
 \|p_Q\|_\infty\le A_h+
       \frac{16\pi}{3}\sqrt{\frac{h}{\varepsilon|\lambda|}}
       \qquad(h>0).
\end{equation}
In particular,
\begin{equation}\label{eq:hub-loglog-density}
 \|p_Q\|_\infty\le\frac12+2\sqrt2+
       \sqrt2\log_+(2B_\varepsilon^2)
                +\frac{16\pi}{3\sqrt{|\lambda|}}.
\end{equation}
For fixed $|\lambda|\ge1$, this is
$O(1+\log\log(1/\varepsilon))$ as $\varepsilon\downarrow0$.
\end{corollary}

\begin{proof}
Apply Lemma~\ref{lem:hub-window-recovery} with the constant in
\eqref{eq:hub-leak-fourier}, and then
Theorem~\ref{thm:hub-loglog-smallball}.
The recovery constant simplifies to
$(2\sqrt{2\pi}/(3\pi))4\sqrt2\pi^{3/2}=16\pi/3$.
\end{proof}

\begin{remark}\label{rem:hub-endpoint-open}
The estimates above hold for each positive $\varepsilon$.
Their constants are not bounded as $\varepsilon\downarrow0$.
Thus they do not establish a uniform three-eigenvalue theorem.
Even though the family at exactly zero leakage is bounded in several
important normal forms, continuity for fixed input laws does not supply
uniformity over all bounded product densities.
\end{remark}

\subsection{Localization and regularity of the hub law}

\begin{proposition}[A localized leakage direction]\label{prop:hub-localized}
Consider
\[
 Q=aX_0^2+2X_0(L+\beta)+\lambda V^2
                          +2\varepsilon LP+\sum_i\phi_i(X_i).
\]
Suppose the coefficient vector of $P$ is supported on $T$, let
$I=T^c$, and put
$G_I=\sum_{i\in I}(w_i,v_i)^T(w_i,v_i)$.
If $G_I$ is invertible and
$(w_i,v_i)G_I^{-1}(w_i,v_i)^T\le1/2$ on $I$, then
\[
 |\mathbb E e^{itQ}|\le
 4\pi^{3/2}(|\lambda|\det G_I)^{-1/2}|t|^{-3/2},
 \qquad
 \|p_Q\|_\infty\le
 3\,2^{4/3}(|\lambda|\det G_I)^{-1/3}.
\]
The estimate is uniform in $\varepsilon$, $a$, and all other displayed
coefficients.
\end{proposition}

\begin{proof}
For fixed $t$, integrating the hub and using $f_0\le1$ bounds the
characteristic function by $(2|t|)^{-1}\int|\widehat n(z)|\,dz$,
where $\widehat n(z)$ is the light expectation with the additional
phase $zL$.  Fresnel linearization of $\lambda V^2$ reduces its
integral to $(4\pi|t\lambda|)^{-1/2}J$, where
\[
 J=\iint\left|\mathbb E e^{izL+iuV+2it\varepsilon LP
                                     +it\sum_i\phi_i(X_i)}\right|dz\,du.
\]
Disintegrate the complex light measure over $P=p$.  On that slice the
cross term translates $z$ by $2t\varepsilon p$.  The triangle
inequality followed by this translation bounds $J$ by the averaged
conditional frame integral.  The coordinates in $I$ are independent
of those carrying $P$; the latter push forward to a complex measure of
total variation at most one.  Therefore
\[
 J\le\iint\prod_{i\in I}
        |\widehat g_i(zw_i+uv_i+tb_i)|\,dz\,du
 \le(4\pi)^2(\det G_I)^{-1/2}.
\]
The last inequality follows by the change of variables associated with
$G_I$ and Lemma~\ref{lem:hub-frame-integral}.
Multiplication of the three factors proves the characteristic-function
bound; its integral proves the density estimate.
\end{proof}

\begin{proposition}[A square on independent coordinates]
\label{prop:hub-free-square}
Let $I$ be a set of light coordinates on which $w$ vanishes and let
$X_R$ denote the others.  Suppose
\[
 Q=2X_0(L+\beta)+\lambda V^2+2Vg(X_R)+F(X_R),
 \qquad \omega=|v_I|>0,
\]
where $g,F$ are arbitrary measurable functions.  Then
\[
 \|p_Q\|_\infty\le\frac12+
             \frac{8\sqrt2}{\omega\sqrt{|\lambda|}}.
\]
\end{proposition}

\begin{proof}
Write $V=V_I+V_R$.  The variable $V_I$ is independent of $X_R$ and
has density at most $\sqrt2/\omega$ by
\eqref{eq:hub-projection}.  Given $X_R$, the remainder is
$\lambda(V_I+c_1)^2+c_2$.  For every Borel set $S$,
\[
 |\{u:\lambda u^2\in S\}|\le2\sqrt{|S|/|\lambda|},
\]
which follows by placing a set of fixed measure at the origin in the
one-sided density of the square.  Thus the conditional probability is
at most $2\sqrt2\omega^{-1}\sqrt{|S|/|\lambda|}$.
Since $L$ is a function of $X_R$ and $\rho_L\le\sqrt2$, its
unnormalized slices have profile
$4\omega^{-1}\sqrt{|S|/|\lambda|}$.
Lemma~\ref{lem:hub-profile} finishes the proof.
\end{proof}

\begin{proposition}[A block carrying only the hub partner]
\label{prop:hub-free-partner}
Assume $a=0$.  Suppose a set $I$ of light coordinates satisfies
$v_i=d_i=0$, $q_i=\kappa w_i$, and $b_i=\mu w_i$ on $I$.
Put $\omega=|w_I|>0$.  Then
\begin{align*}
 \|p_Q\|_\infty
 &\le\frac12+\frac{\sqrt2}{\omega}
 \left[\frac72\left(2+\log(1+4\pi^4/|\lambda|)\right)
                          +\frac12\log(1+4|\kappa|)\right].
\end{align*}
If $\kappa=0$, the sharper bound is
\[
 \|p_Q\|_\infty\le\frac12+\frac{\sqrt2}{\omega}
                        \left(2+\log(1+2\pi^4/|\lambda|)\right).
\]
\end{proposition}

\begin{proof}
Let $R=I^c$, $L_I=\langle w_I,X_I\rangle$ and $L_R=L-L_I$.
Collecting terms gives
$Q=2X_0t+2\kappa L^2+2LP+W$, where $t=L+\beta$ and $P,W$
depend only on $X_R$.  The value of the hub at level $s$ has the form
\[
 x_* =\frac{\alpha}{2t}+c_0-\kappa t,
 \quad
 \alpha=s-2\kappa\beta^2+2\beta P-W,
 \quad c_0=2\kappa\beta-P.
\]
Conditional on $X_R$, the variable $t$ has density at most
$\sqrt2/\omega$.  On $|t|\ge1$ its contribution to $p_Q$ is at most
$1/2$.  Thus the remaining contribution is bounded by
\[
 \frac{\sqrt2}{\omega}\mathbb E_{X_R}
 \int_{0<|t|<1}\widetilde f_0\!\left(\frac\alpha{2t}-\kappa t\right)
                                          \frac{dt}{2|t|},
 \qquad\widetilde f_0(x)=f_0(x+c_0).
\]
For $\kappa=0$, the substitution $x=\alpha/(2t)$ and the bathtub
principle bound the inner integral by $\log(1+1/|\alpha|)$.
For $\kappa\ne0$, set $\sigma=2\kappa$ and
$\tau=\min(1,\sqrt{|\alpha|/(2|\sigma|)})$.
For $0<|t|<\tau$, the map
$x(t)=\alpha/(2t)-\sigma t/2$ is monotone on each half-line,
its two images are disjoint, $|x(t)|\ge|\alpha|/(4\tau)$, and
$dt/|t|\le3\,|dx|/|x|$.
The bathtub principle bounds this part by
$3\log(1+2\tau/|\alpha|)$; the range $\tau\le|t|<1$ contributes at
most $\log(1/\tau)$.
Together these are at most
\[
 \frac72\log(1+2/|\alpha|)+\frac12\log(1+2|\sigma|).
\]
The variable $\alpha$ is a quadratic polynomial on $X_R$ consisting of
$-\lambda V^2$ and one-coordinate terms: the $L$-cross terms have
already been put into $P$.  One-dimensional Fresnel linearization and
Lemma~\ref{lem:hub-frame-integral} give
$|\mathbb E e^{it\alpha}|\le2\sqrt{\pi/|t\lambda|}$.
The kernel in Lemma~\ref{lem:hub-frozen} therefore gives
$\mathbb P(|\alpha-u|\le r)\le\pi^2\sqrt{2r/|\lambda|}$.
For $c=2\pi^2/\sqrt{|\lambda|}$, layer cake implies
\[
 \mathbb E\log(1+2/|\alpha|)
 \le\int_0^\infty\min\{1,c(e^u-1)^{-1/2}\}\,du
 \le\log(1+c^2)+2.
\]
Indeed, above $u=\log(1+c^2)$ the integral equals
$2c\arctan(1/c)<2$.
Replacing $2$ by $1$ in the logarithm gives
$2+\log(1+2\pi^4/|\lambda|)$, proving both assertions.
\end{proof}

\begin{proposition}[A hub density of bounded variation]
\label{prop:hub-bv}
Assume $a=0$ and let the total variation of $f_0$ be $V_0<\infty$.
For $A=4\sqrt{\pi/|\lambda|}$,
\begin{equation}\label{eq:hub-bv-fourier}
 |\mathbb E e^{itQ}|\le V_0|t|^{-3/2}
       \left[A\left(1+\log_+(2|t|^{3/2}/V_0)\right)+\frac12\right].
\end{equation}
In particular, for fixed $|\lambda|\ge1$,
$\|p_Q\|_\infty\le C V_0^{2/3}$ with an absolute constant $C$,
uniformly in every light diagonal, linear term, and cross term.
\end{proposition}

\begin{proof}
Integration by parts against the distributional derivative of $f_0$
gives $|\widehat f_0(z)|\le\min(1,V_0/|z|)$.
Freeze $L$ and apply \eqref{eq:hub-slice-fourier}.  With $u=L+\beta$,
\[
 |\mathbb E e^{itQ}|\le
 \int\min\{1,V_0/(2|tu|)\}
       \min\{\rho_L(u-\beta),A|t|^{-1/2}\}\,du.
\]
On $|u|\le T=|t|^{1/2}$, use the second entry in the second minimum.
The integral of the first minimum is at most
$(V_0/|t|)(1+\log_+(2|t|T/V_0))$.
On $|u|>T$, retain $\rho_L$ and use its unit mass, obtaining
$V_0/(2|t|^{3/2})$.  This proves
\eqref{eq:hub-bv-fourier}.
Substitute $t=V_0^{2/3}u$ into the integral of the minimum of one and
the right side.  The resulting integral is finite and uniformly
bounded for $A\le4\sqrt\pi$, proving the density estimate.
No numerical quadrature is needed for the existence of $C$.
\end{proof}

\subsection{Bounded light supports and a fixed transverse projection}

\begin{proposition}[A maximal-density criterion]
\label{prop:hub-maximal-projection}
Let $E$ be a fixed $k$-dimensional subspace orthogonal to $w,v$,
and write $Z$ for the orthonormal coordinates of the projection of
$X'$ onto $E$.  Let
\[
 Q=2X_0(L+\beta)+\lambda V^2+2Vg(L,Z)+F(L,Z),
\]
where $g,F$ are measurable.  Put
\[
 M_k=\mathop{\rm ess\,sup}_{\ell}
            \int_{\mathbb R^k}\mathop{\rm ess\,sup}_u
                       \rho_{L,V,Z}(\ell,u,z)\,dz.
\]
If $M_k<\infty$, then
\[
 \|p_Q\|_\infty\le\frac12+8M_k|\lambda|^{-1/2}.
\]
This result needs no leverage condition.
\end{proposition}

\begin{proof}
On the slice $L=\ell$ complete the square in $V$:
$\lambda(V-V_0(z))^2+F_0(z)$.
The slab $|V-V_0(z)|\le r$ has slice mass at most $2rM_k$.
Outside the slab, the two inverse branches each have derivative
at most $(2|\lambda|r)^{-1}$ in the output variable, so the
restricted output density is at most $M_k/(|\lambda|r)$.
For every Borel set $S$, choose
$r=\sqrt{|S|/(2|\lambda|)}$ to obtain the profile
$2\sqrt2M_k\sqrt{|S|/|\lambda|}$.
Lemma~\ref{lem:hub-profile} gives the result; freezing is unnecessary
here because the remainder is already a function of $(L,V,Z)$.
\end{proof}

\begin{proposition}[Bounded light supports]\label{prop:hub-light-intervals}
Suppose every light coordinate is supported in an interval of length
at most $K$.  Then the quantity in
Proposition~\ref{prop:hub-maximal-projection} satisfies
\[
 M_k\le2^{(k+2)/2}K^k
  \int_{\mathbb R^k}e^{-2(|y|-k-2)_+^2}\,dy.
\]
Moreover, for \eqref{eq:hub-normal} with $a=d_i=0$, arbitrary $q$
and arbitrary linear terms,
\begin{equation}\label{eq:hub-light-interval-bound}
 \|p_Q\|_\infty\le\frac12+
 8\sqrt2\left[2^{3/2}\left(6+\sqrt{\pi/2}\right)
                         \frac K{|\lambda|}\right]^{1/2}
 <\frac12+52\sqrt{K/|\lambda|}.
\end{equation}
The hub law is arbitrary subject to its density cap.
\end{proposition}

\begin{proof}
Tilt the product law by $e^{\langle\theta,Z\rangle}$.
Each tilted coordinate density is at most $e^{|\theta|K}$ times
its original cap, because the tilting factor varies by at most
that amount on its supporting interval.  The $(k+2)$-dimensional
projection-density bound for the tilted product law is therefore
$2^{(k+2)/2}e^{(k+2)|\theta|K}$.
Independence and the bounded-variable estimate of
\citet{Hoeffding1963} give
\[
 \mathbb E e^{\langle\theta,Z-\mathbb EZ\rangle}
                 \le e^{|\theta|^2K^2/8}.
\]
Undoing the tilt and optimizing $\theta$ in the direction
$z-\mathbb EZ$ yields
\[
 \rho_{L,V,Z}(\ell,u,z)
 \le2^{(k+2)/2}
  \exp\left[-2\frac{(|z-\mathbb EZ|-(k+2)K)_+^2}{K^2}\right].
\]
The optimization uses
$|\theta|=4(|z-\mathbb EZ|-(k+2)K)_+/K^2$.
Integration over $z$ proves the first assertion.

For the second assertion freeze $L$ as in
Remark~\ref{rem:hub-freezing}.  With no light diagonal, every frozen
remainder has the form
$\lambda(V-V_0)^2+\gamma B+c$, where $B$ is a unit linear
projection orthogonal to $w,v$; its direction may depend on the
frozen parameter.  The preceding tilt bound is uniform in that
direction and gives
\[
 \int\sup_u\rho_{L,V,B}(\ell,u,b)\,db
 \le M:=2^{3/2}\left(6+\sqrt{\pi/2}\right)K.
\]
The slab $|V-V_0|\le r$ has slice mass at most $4r$, using
$\rho_{L,V}\le2$.  Off the slab the output density is at most
$M/(|\lambda|r)$.  Optimizing
$4r+M|S|/(|\lambda|r)$ gives the Borel profile
$4\sqrt{M|S|/|\lambda|}$.
Lemma~\ref{lem:hub-profile} proves
\eqref{eq:hub-light-interval-bound}.  If $\gamma=0$, the direct
square profile is at least as strong, since $K\ge1$.
\end{proof}

\begin{corollary}[Window-mixture light laws]\label{cor:hub-window-lights}
Suppose each light law is of the form $\mathcal U[0,1]*\nu_i$.
For \eqref{eq:hub-normal} with $a=d_i=0$, every independent cap-one
hub law satisfies
\[
 \|p_Q\|_\infty\le\frac12+52|\lambda|^{-1/2}.
\]
\end{corollary}

\begin{proof}
Represent $X_i=U_i+Z_i$, with independent uniform $U_i$ on $[0,1]$,
and condition on the shifts $Z_i$.  The light coordinates are now
supported on unit intervals.  The shifts change only the linear
coefficients and the constant of the quadratic form, which are
unrestricted in Proposition~\ref{prop:hub-light-intervals}.
Apply that proposition and average.
\end{proof}

\subsection{Several hub coordinates}

\begin{theorem}[Three or more hubs]\label{thm:hub-many}
Let $X_H\in\mathbb R^k$ and $X'$ be independent, with
$\int f_H^2\le1$, and suppose $L=MX'\in\mathbb R^k$ has density
bounded by $D$.  For any real measurable $F,R$ and $\beta\in\mathbb R^k$,
\[
 Q=F(X_H)+2\langle X_H,L+\beta\rangle+R(X')
\]
satisfies
\[
 |\mathbb E e^{itQ}|\le
       \min\{1,D^{1/2}(\pi/|t|)^{k/2}\}.
\]
If $k>2$, then
\begin{equation}\label{eq:hub-many-density}
 \|p_Q\|_\infty\le\frac{k}{k-2}D^{1/k}.
\end{equation}
For independent cap-one light coordinates and singular values
$s_1,\ldots,s_k>0$ of $M$, this is at most
$\frac{k}{k-2}\sqrt2(s_1\cdots s_k)^{-1/k}$.
\end{theorem}

\begin{proof}
Let
$h_t(\ell)=\int f_H(x)e^{itF(x)+2it\langle x,\ell+\beta\rangle}\,dx$.
Plancherel gives
\[
 \int_{\mathbb R^k}|h_t(\ell)|^2\,d\ell
   =(\pi/|t|)^k\int f_H^2\le(\pi/|t|)^k.
\]
After conditioning on $X'$,
\[
 |\mathbb E e^{itQ}|\le\mathbb E|h_t(L)|
 \le\|\rho_L\|_2\|h_t\|_2
 \le D^{1/2}(\pi/|t|)^{k/2}.
\]
Integrating the minimum with one gives
\eqref{eq:hub-many-density}.  The final assertion follows from
\eqref{eq:hub-projection}.
\end{proof}

\begin{theorem}[Two hubs and a transverse square]\label{thm:hub-two-square}
Let $X_H\in\mathbb R^2$ have density $f_H\le1$ and be independent
of the cap-one product vector $X'$.  Let $M$ have rank two, singular
values $s_1,s_2$, and orthonormal row-space basis $e_1,e_2$.
Suppose $v\perp e_1,e_2$, $|v|=1$, and
$e_{1i}^2+e_{2i}^2+v_i^2\le1/2$.
Put $L=MX'$ and $V=\langle v,X'\rangle$.  For every real symmetric
$2\times2$ matrix $B$, vectors $q_1,q_2$, and coordinate quadratics,
\[
 Q=X_H^TBX_H+2\langle X_H,L+\beta\rangle+\lambda V^2
                   +2\sum_{j=1}^2L_j\langle q_j,X'\rangle
                   +\sum_i\phi_i(X_i)
\]
has
\[
 \|p_Q\|_\infty\le
 5\,2^{6/5}(s_1s_2)^{-2/5}|\lambda|^{-1/5}
 <11.5(s_1s_2)^{-2/5}|\lambda|^{-1/5}.
\]
\end{theorem}

\begin{proof}
Freeze the two partners on rectangular cells.  Plancherel gives
\[
 \int|h_t(\ell)|^2\,d^2\ell\le\pi^2/t^2
\]
for the hub factor.
Fresnel linearization of the transverse square and the
three-dimensional frame bound give for the frozen slice
\[
 |\widehat\mu_{\ell_0,\ell}(t)|
 \le\min\left\{\rho_L(\ell),
            \frac{8\sqrt\pi}{s_1s_2\sqrt{|t\lambda|}}\right\}.
\]
The factor $(s_1s_2)^{-1}$ is the Jacobian converting the frequency
of $L$ to that of its orthonormal row-space coordinates.
For each fixed $t$, Cauchy--Schwarz on the cells and passage to the
Riemann-sum limit give
\[
 |\mathbb E e^{itQ}|\le K|t|^{-5/4},
 \quad K=\pi\left(\frac{8\sqrt\pi}{s_1s_2}\right)^{1/2}
                              |\lambda|^{-1/4}.
\]
Indeed, the squared slice factor integrates to at most its displayed
uniform bound times $\int\rho_L=1$.
The characteristic function is also at most one.  Its integral is
$10K^{4/5}$, and division by $2\pi$ gives the stated constant.
\end{proof}

\begin{theorem}[Two hubs with an interval slice profile]
\label{thm:hub-two-profile}
Let $X_H\in\mathbb R^2$ have density at most one, be independent of
$X'$, and let $L=MX'$ have a density.  Let $Y_\ell$ be a family of
real functions of $X'$ such that $Y_{\ell'}\to Y_\ell$ pointwise as
$\ell'\to\ell$.  Suppose its frozen slices satisfy
\[
 \mu_{\ell_0,\ell}(I)\le K|I|^b
 \quad(0<b\le1)
\]
for every fixed $\ell_0$, almost every $\ell$, and every interval $I$.
Then, for every real symmetric $B$ and $\beta\in\mathbb R^2$,
\[
 Q=X_H^TBX_H+2\langle X_H,L+\beta\rangle+Y_L
\]
has
\[
 \|p_Q\|_\infty\le\frac1\pi+
             \frac{\sqrt{10K}}{1-2^{-b/2}}.
\]
\end{theorem}

\begin{proof}
First restrict the light measure to $\{\rho_L(L)\le N\}$.
The resulting subprobability measure has bounded partner density, and
all the assumed slice inequalities persist with the same constants.
The bounds below are independent of $N$; monotone convergence at the
end removes the restriction.  Use rectangular cells $C_j$ and frozen
values $\ell_j$.  Define the
finite positive measure
\[
 \nu_h=\sum_j\int_{C_j}
          \mu_{\ell_j,\ell}*\widetilde\mu_{\ell_j,\ell}\,d^2\ell.
\]
It has nonnegative Fourier transform
$\widehat\nu_h(t)=\sum_j\int_{C_j}|\widehat\mu_{\ell_j,\ell}(t)|^2
d^2\ell$ and satisfies
$\nu_h([-r,r])\le K(2r)^b$, since each slice has mass
$\rho_L(\ell)$ and $\int\rho_L=1$.

For any positive measure $\nu$ with nonnegative transform and
$\nu([-r,r])\le K'r^b$, the Fej\'er kernel gives
\[
 \int_{-T}^T\widehat\nu(t)\,dt
 \le4T\int\left(\frac{\sin(Tx)}{Tx}\right)^2\nu(dx)
 \le20K'T^{1-b}.
\]
For the last inequality split at $|x|=1/T$ and into the dyadic shells
$2^j/T<|x|\le2^{j+1}/T$; use
$\min(1,(Tx)^{-2})$ and
$1+2^b/(1-2^{b-2})\le5$.

Let $\Omega_A=\{A\le|t|\le2A\}$.  Cauchy--Schwarz in the partners
and then in $t$, followed by Fatou and the freezing limit, give
\[
 \int_{\Omega_A}|\mathbb E e^{itQ}|\,dt
 \le \left(\int_{\Omega_A}\int|h_t(\ell)|^2d^2\ell\,dt\right)^{1/2}
       \liminf_{h\downarrow0}
         \left(\int_{\Omega_A}\widehat\nu_h(t)dt\right)^{1/2}.
\]
This can first be done on finitely many cells and then exhausted; all
slice-mass estimates remain unchanged.  The first factor is at most
$\pi A^{-1/2}$ by hub Plancherel, and the second is at most
$(20\cdot2^bK(2A)^{1-b})^{1/2}$.
Thus the integral is at most $\pi\sqrt{40K}A^{-b/2}$.
Sum over $A=2^j$, $j\ge0$, add the bound two on $|t|<1$, and divide
by $2\pi$.
\end{proof}

\begin{proposition}[A projected crossing remainder]
\label{prop:hub-crossing}
Let $k=1$ or $2$, let $X_H\in\mathbb R^k$ have density at most one,
and let $X'$ be an independent cap-one product vector.  Consider
\[
 Q=X_H^TBX_H+2\langle X_H,MX'+\beta\rangle+R(X'),
 \qquad R(x)=x^TA_Rx+b^Tx+c.
\]
Split the light coordinates as $I\cup J$.  Suppose $M_J$ has rank
$k$, with product $S$ of its singular values, and put
\[
 C=(A_R)_{I,J}(I-\Pi),\qquad F=\|C\|_F>0,
\]
where $\Pi$ projects onto the row space of $M_J$.
If every row $c_i$ of $C$ satisfies $|c_i|^2\le F^2/2$, then
\begin{equation}\label{eq:hub-crossing-fourier}
 |\mathbb E e^{itQ}|\le
 2^{(k+3)/4}\pi^{(k+1)/2}(SF)^{-1/2}|t|^{-(k+1)/2}.
\end{equation}
For $k=2$ this gives
\[
 \|p_Q\|_\infty\le3\,2^{5/6}(SF)^{-1/3}<5.35(SF)^{-1/3}.
\]
For $k=1$, $S=|w_J|$, it gives
\[
 \mathbb P(|Q-s|\le\delta)\le
 \frac{\pi^2\delta}{\sqrt{SF}}
 \left(1+\log_+\frac{\sqrt{SF}}{4\delta}\right).
\]
\end{proposition}

\begin{proof}
Let $\mu_t(\ell)$ be the density of
$\mathbb E[e^{itR(X')};MX'\in d\ell]$.
The hub factor and Plancherel imply
\[
 |\mathbb E e^{itQ}|^2\le(\pi/|t|)^k\int|\mu_t(\ell)|^2d\ell.
\]
Convolve $\mu_t$ with an even nonnegative normalized approximate
identity $k_h$.  Conditional Cauchy--Schwarz in $X_I$, an independent
copy $X_J^*$, and integration over $\ell$ yield
\[
 \int|\mu_t*k_h|^2\le
 \mathbb E\left[\kappa_h(M_JD_J)
       \prod_{i\in I}|\widehat f_i(2t\langle a_i,D_J\rangle)|\right],
\]
where $D_J=X_J-X_J^*$, $\kappa_h=k_h*k_h$, and $a_i$ is the
$i$th row of $(A_R)_{I,J}$.  Terms of $R$ internal to $I$ cancel;
terms internal to $J$ contribute only a phase of modulus one.

Put $\theta_i=|c_i|^2/F^2$ and discard the zero rows, whose factors
are at most one.  H\"older with weights $\theta_i$ bounds the last
expectation by
\[
 \prod_i\left(\mathbb E\left[\kappa_h(M_JD_J)
 |\widehat f_i(2t\langle a_i,D_J\rangle)|^{1/\theta_i}\right]
                                                   \right)^{\theta_i}.
\]
The coordinates of $D_J$ are independent with densities bounded by
one.  Since $c_i\perp\operatorname{row}(M_J)$, the density of
$(M_JD_J,\langle c_i,D_J\rangle)$ is at most
$2^{(k+1)/2}/(S|c_i|)$.
Write $a_i=c_i+\alpha_iM_J$ and translate the last variable for each
value of $M_JD_J$.  Applying \eqref{eq:hub-scalar-lp} gives a bound
\[
 \frac{2^{(k+1)/2}}{S|c_i|}\frac{4\pi\sqrt{\theta_i}}{2|t|}
 =\frac{2^{(k+3)/2}\pi}{SF|t|}
\]
for every factor.  It is independent of $h$.
Pass to the $L^2$ approximate-identity limit and combine with the
hub factor to obtain \eqref{eq:hub-crossing-fourier}.
For $k=2$, integrate its minimum with one.
For $k=1$, put $K_1=2\pi/\sqrt{SF}$ and apply the kernel used in
Lemma~\ref{lem:hub-frozen} to
$\min(1,K_1/|t|)$.  The integral over
$|t|\le\pi/(2\delta)$ is at most
$2K_1(1+\log_+(\pi/(2K_1\delta)))$, which gives the assertion.
\end{proof}

\begin{remark}
The projected-crossing estimate permits arbitrary terms internal to
either light block.  Its extension to several volume directions and
arbitrary-rank crossing matrices is given in
Theorem~\ref{thm:grouped-crossing}; exact rank of the crossing matrix
is not needed there.  The nondegeneracy and volume-marginal assumptions
remain essential hypotheses of that theorem.
\end{remark}

\begin{corollary}[A logarithmic bound with a hub diagonal]
\label{cor:hub-diagonal-log}
Suppose Theorem~\ref{thm:hub-leak} and the $k=1$ case of
Proposition~\ref{prop:hub-crossing} both apply.  Then
\[
 \|p_Q\|_\infty\le
 \frac{2}{\sqrt{SF}}
 \left(3+\log_+\frac{4(SF)^{3/2}}{\varepsilon|\lambda|}\right).
\]
\end{corollary}

\begin{proof}
The characteristic function is bounded by
$\min(1,K_1|t|^{-1},K_2|t|^{-3/2})$, where
$K_1=2\pi/\sqrt{SF}$ and
$K_2=4\sqrt2\pi^{3/2}/\sqrt{\varepsilon|\lambda|}$.
Splitting the positive half-line at $K_1$ and $(K_2/K_1)^2$ gives
an integral at most
$K_1[3+\log_+(K_2^2/K_1^3)]$.
The same bound holds when the latter split precedes the former by
using just the first and third entries of the minimum.
Since $K_2^2/K_1^3=4(SF)^{3/2}/(\varepsilon|\lambda|)$,
division of the two-sided integral by $2\pi$ proves the claim.
\end{proof}

\begin{proposition}[A spread linear hub]\label{prop:hub-spread}
Let $X_H$ and $X'$ be independent cap-one product vectors, let
$U=\langle c,X_H\rangle$, and assume
$c_h^2\le|c|^2/2$ for every hub coordinate.  For
\[
 Q=2U(L+\beta)+\sum_{h\in H}\phi_h(X_h)
       +\lambda V^2+2L\langle q,X'\rangle+\sum_i\phi_i(X_i),
\]
with \eqref{eq:hub-frame},
\[
 |\mathbb E e^{itQ}|\le
 8\pi^{3/2}|c|^{-1}|\lambda|^{-1/2}|t|^{-3/2},
 \qquad
 \|p_Q\|_\infty\le12(|c|^2|\lambda|)^{-1/3}.
\]
\end{proposition}

\begin{proof}
The conditional hub expectation is
\[
 H_t(u)=\prod_{h\in H}
  \widehat g_h(2tuc_h+tb_h),\qquad
 g_h(x)=f_h(x)e^{itd_hx^2}.
\]
H\"older with exponents $p_h=|c|^2/c_h^2\ge2$ and
\eqref{eq:hub-scalar-lp} gives
$\int|H_t(u)|\,du\le2\pi/(|t||c|)$.
Freeze $L$ and pair this factor with the uniform slice-transform bound
$4\sqrt\pi/\sqrt{|t\lambda|}$ from
Lemma~\ref{lem:hub-frozen}.  Their product is the claimed
characteristic-function bound.  Integrating the minimum with one gives
the density bound.
\end{proof}

\subsection{An adapted-hub variational formulation}

The entangled case admits a variational formulation in terms of a
restricted conditional law.  Its dependence structure is absent from
arbitrary bounded two-dimensional densities.

\begin{proposition}[Adapted-hub variational reduction]
\label{prop:hub-adapted}
Let $a=0$ and $q=q_0+\varepsilon e$, where
$q_0\in\operatorname{span}\{w,v\}$ and $P=\langle e,X'\rangle$.
Absorb $-2\varepsilon\beta P$ into the light linear terms and write
\[
 Q=2Y(L+\beta)+R_0(X'),\qquad Y=X_0+\varepsilon P.
\]
Then the conditional density of $Y$ given $X'$ is at most one and
depends on $X'$ only through $P$.
For a fixed level $s$, put
$x_*=(s-R_0(X'))/(2(L+\beta))$.
If $W_s(x,p)$ is the joint density of $(x_*,P)$ under the
nonnegative weight $1/(2|L+\beta|)$, then the supremum at $s$ over
all conditional hub densities $f(y\mid p)\le1$ is
\begin{equation}\label{eq:hub-adapted-functional}
 \Phi(s)=\int_{\mathbb R}
       \sup_{|A|=1}\int_A W_s(x,p)\,dx\,dp.
\end{equation}
In particular, $p_Q(s)\le\Phi(s)$ for the independent-hub form.
The same assertion can be formulated for window probabilities, without
assuming that the weighted pushforward has a density.
\end{proposition}

\begin{proof}
The identity
$2X_0(L+\beta)+2\varepsilon LP
 =2(L+\beta)(X_0+\varepsilon P)-2\varepsilon\beta P$
proves the first assertion.  Its conditional density is
$f_0(y-\varepsilon P)$.
Given $X'$, the map from $Y$ to $Q$ is affine with derivative
$2(L+\beta)$; hence
\[
 p_Q(s)=\iint f(x\mid p)W_s(x,p)\,dx\,dp.
\]
For each $p$, maximizing a nonnegative linear functional subject to
$0\le f(\cdot\mid p)\le1$ and $\int f(\cdot\mid p)=1$
gives the integral over a set of measure one on which $W_s$ is largest.
Measurable approximate maximizers can be chosen by thresholding
$W_s$ and filling a level set; monotone approximation then gives
\eqref{eq:hub-adapted-functional}.  The window version follows by
applying the same conditional maximization to the nonnegative
conditional window kernel before taking the density limit.
\end{proof}

\begin{conjecture}[A sufficient adapted-hub estimate]
\label{conj:hub-adapted}
Let $X'$ have independent coordinate densities bounded by one, and let
$w,v,e$ be orthonormal with $w_i^2+v_i^2+e_i^2\le1/2$.
Set $L=\langle w,X'\rangle$, $V=\langle v,X'\rangle$, and
$P=\langle e,X'\rangle$.
If the conditional density of $Y$ given $X'$ is $g(y\mid P)\le1$
of unit mass, then $2YL+V^2$ has density bounded by an absolute constant.
Equivalently, its functional \eqref{eq:hub-adapted-functional} is
bounded uniformly in the level.
\end{conjecture}

The conjecture is stronger than the estimate for a translated
independent hub, because it permits an arbitrary probability density
depending on $P$.  This conditional estimate remains open.  The
preceding estimates also do not cover arbitrary quadratic twists on
the orthogonal complement of $\operatorname{span}\{w,v\}$ or yield an
exhaustive quantitative reduction of all three-eigenvalue matrices
to the stated hub configurations.

\subsection{Adapted hubs for Gaussian mixtures}

The adapted-hub estimate can be proved for a class allowing arbitrarily
small Gaussian component variances.  The resulting theorem concerns the
entangled core $2YL+\lambda V^2$; it does not assert an estimate for
arbitrary light diagonal terms.

\begin{lemma}[A conditional Gaussian bound]\label{lem:hub-gaussian}
Suppose $(L,V,P)$ is a Gaussian vector with covariance matrix at least
$hI_3$, where $h>0$.  Let the conditional law of $Y$ given $(L,V,P)$
have density $g(y\mid P)\le1$ and unit integral in $y$.
For $|\lambda|\ge1$ and arbitrary finite real measurable $B,C$, set
\[
 Q=2YL+\lambda V^2+B(Y,P)V+C(Y,P).
\]
Then
\begin{equation}\label{eq:hub-gaussian-bound}
 \|p_Q\|_\infty\le4(2\pi)^{-3/4}h^{-3/4}.
\end{equation}
If $h\ge(2\pi)^{-1}$, the sharper numerical bound is
$\|p_Q\|_\infty\le3$.
The means of the Gaussian vector are arbitrary.
\end{lemma}

\begin{proof}
Condition on $P$.  The conditional covariance of $(L,V)$ is at least
$hI_2$.  Indeed, for the full covariance matrix $M\ge hI_3$, its
Schur complement satisfies
\[
 a^TK a=\inf_b(a,b)^TM(a,b)\ge h|a|^2.
\]
Regress $L$ on $V$: $L=\mu+\alpha V+W$, where $W$ is an independent
Gaussian with variance at least $h$, and $V$ also has variance at
least $h$.  Given $P$, the variable $Y$ is independent of these
Gaussian residuals.  For fixed $Y=y$, the $W$ factor and the
one-dimensional quadratic Gaussian integral imply
\[
 |\mathbb E[e^{itQ}\mid P,Y=y]|
 \le e^{-2hy^2t^2}(1+4h^2t^2)^{-1/4}.
\]
The added terms change only a linear phase in $V$ and a constant.
Completing the Gaussian square shows that the modulus contributed by
the linear phase is at most one, for every conditional mean.
Using both $g\le1$ and $\int g=1$ gives
\begin{equation}\label{eq:hub-gaussian-envelope}
 |\mathbb E e^{itQ}|\le
 (1+4h^2t^2)^{-1/4}
       \min\left\{1,\frac{\sqrt\pi}{\sqrt{2h}|t|}\right\}.
\end{equation}
This integrable bound is independent of $P$.
Set $a=\sqrt{2\pi h}$ and $u=2h|t|$.  Fourier inversion gives
\[
 \|p_Q\|_\infty\le a^{-2}\int_0^\infty
             (1+u^2)^{-1/4}\min(1,a/u)\,du.
\]
Bound $(1+u^2)^{-1/4}$ by $u^{-1/2}$ and split at $u=a$.
The integral is at most $4\sqrt a$, proving
\eqref{eq:hub-gaussian-bound}.
For $h=(2\pi)^{-1}$, $a=1$; splitting at one and using one below
that point and $u^{-1/2}$ above it gives the sharper value $1+2=3$.
The envelope decreases with $h$, so this holds for larger $h$ too.
\end{proof}

\begin{corollary}[A Gaussian covariance floor]\label{cor:hub-gaussian-floor}
Let $X'=Z+G$, where $Z$ is an arbitrary random vector and $G$ is an
independent Gaussian vector with covariance at least $(2\pi)^{-1}I$.
Let $L,V,P$ be orthonormal linear projections and let $Y$ have a
conditional density at most one depending on $X'$ only through $P$.
Then every $Q$ in Lemma~\ref{lem:hub-gaussian} has density at most
three.  No spread hypothesis on the frame is needed.
\end{corollary}

\begin{proof}
Generate the conditional hub after $X'$, using independent auxiliary
randomness.  Conditional on $Z$, the vector $(L,V,P)$ is Gaussian with
the required covariance floor and with arbitrary mean.  Its hub kernel
is still $g(y\mid P)$.  Apply Lemma~\ref{lem:hub-gaussian} and average.
\end{proof}

\begin{lemma}[A uniform partition of fixed-dimensional frames]
\label{lem:hub-frame-partition}
For each positive integer $d$ there is $\eta_d>0$ such that every
finite family $r_i\in\mathbb R^d$ satisfying
\[
 \sum_i r_i r_i^T=I_d,\qquad |r_i|^2\le\tfrac12
\]
can be partitioned into $A,B$ with
\begin{equation}\label{eq:hub-frame-partition}
 \sum_{i\in A}r_i r_i^T\ge\eta_d I_d,
 \qquad
 \sum_{i\in B}r_i r_i^T\ge\eta_d I_d.
\end{equation}
The constant does not depend on the number of vectors.
\end{lemma}

\begin{proof}
We first recall the finite vector-matroid base-packing criterion:
a family of rank $d$ contains two disjoint bases if and only if
$|E\setminus F|\ge2(d-\operatorname{rank}F)$ for every subfamily $F$
\citep{Edmonds1965}.  In a Parseval frame as above, the energy
orthogonal to a subspace of codimension $k$ is $k$.
Since a single vector contributes at most $1/2$, at least $2k$
vectors lie outside that subspace.  Thus two disjoint bases exist
for each fixed frame.

Suppose the uniform assertion were false.  Choose frames whose best
lower bound in \eqref{eq:hub-frame-partition} tends to zero.
Sort each frame in decreasing order of vector norm, pad with zero
vectors, and take a diagonal subsequence so that its $i$th vector
converges to $r_i$ for each fixed $i$.
Because the total squared norm is $d$, every tail vector beyond
index $N$ has squared norm at most $d/(N+1)$.
The limiting atoms satisfy $\sum_i r_i r_i^T\le I_d$.
Put
\[
 S=I_d-\sum_i r_i r_i^T\ge0,\qquad K=\ker S.
\]
The projections of the limiting atoms onto $K$ form a Parseval frame
for $K$, each with squared norm at most $1/2$.
They contain two disjoint finite bases of $K$.
To justify this for the possibly infinite family, choose a finite
initial segment whose covariance on $K$ is at least
$(1-\epsilon)I_K$, with $\epsilon<1/(2d)$.
Every codimension-$k$ subspace of $K$ then has outside energy greater
than $k-1/2$, so the number of outside vectors is at least $2k$.
The finite base-packing criterion applies to this initial segment.

Let $A_0,B_0$ be the indices of these disjoint bases, taking both
empty if $K=\{0\}$.  The matrices
\[
 H_A=\sum_{i\in A_0}r_i r_i^T+S/2,
 \qquad
 H_B=\sum_{i\in B_0}r_i r_i^T+S/2
\]
are positive definite.  A null vector would belong to $K$ and be
orthogonal to a basis of $K$.  Let $\gamma>0$ be the smaller of their
least eigenvalues.  Choose $N$ containing both bases and large enough
that $d/(2\sqrt{N+1})<\gamma/4$.

In each approximating frame assign the indicated basis vectors to
their prescribed sides, and assign the other first $N$ vectors
arbitrarily.  Partition the tail independently by fair coin flips.
If $T$ is the tail covariance and $T_A$ its contribution to side $A$,
\[
 \mathbb E\|T_A-T/2\|_F^2
 =\tfrac14\sum_{i>N}|r_i^{\rm app}|^4
 \le\frac{d^2}{4(N+1)}.
\]
There is consequently a tail partition with operator error at most
$d/(2\sqrt{N+1})$ on both sides.  For fixed $N$, the tail covariance
converges to $I_d-\sum_{i\le N}r_i r_i^T\ge S$.
The prescribed basis covariances converge as well.  Thus sufficiently
late approximating frames have both partition covariances at least
$\gamma I_d/2$, a contradiction.
\end{proof}

\begin{theorem}[An adapted hub for Gaussian component mixtures]
\label{thm:hub-gaussian-mixtures}
Let
\[
 X_i=M_i+\sqrt{T_i}\,G_i,\qquad T_i>0,
\]
where the pairs $(M_i,T_i)$ are independent across coordinates and
the standard Gaussian variables $G_i$ are independent of all these
pairs.  Suppose
\begin{equation}\label{eq:hub-mixture-peak}
 \mathbb E T_i^{-1/2}\le\sqrt{2\pi}\qquad\text{for every }i.
\end{equation}
Let $w,v,e$ be orthonormal with
$w_i^2+v_i^2+e_i^2\le1/2$, put
$(L,V,P)=(\langle w,X'\rangle,\langle v,X'\rangle,
\langle e,X'\rangle)$, and let $Y$ have conditional density
$g(y\mid P)\le1$ of unit mass.
For $|\lambda|\ge1$ and arbitrary finite real measurable $B,C$,
\[
 Q=2YL+\lambda V^2+B(Y,P)V+C(Y,P)
\]
has density bounded by
\begin{equation}\label{eq:hub-mixture-bound}
 4\,2^{-3/4}3^{3/2}\eta_3^{-3/4},
\end{equation}
where $\eta_3$ is the constant in
Lemma~\ref{lem:hub-frame-partition}.
In particular this gives a uniform bound for
\[
 aX_0^2+2X_0(L+\beta)+\lambda V^2+2\varepsilon LP
\]
for every independent cap-one hub law, every $a,\beta,\varepsilon$,
and every light law satisfying \eqref{eq:hub-mixture-peak}.
\end{theorem}

\begin{proof}
Put $r_i=(w_i,v_i,e_i)^T$ and choose a partition $A,B$ from
Lemma~\ref{lem:hub-frame-partition}.  Its two covariance sums satisfy
$S_A\ge\eta_3I_3$ and $S_B\ge\eta_3I_3$.
Realize $Y$ after $X'$ using fresh auxiliary randomness and the
prescribed kernel $g(y\mid P)$.  This preserves the joint law relevant
to $Q$ and ensures that the same kernel remains valid after
conditioning on the latent mixture variables.
Conditional on the mixture variables, the covariance of $(L,V,P)$ is
\[
 M=\sum_iT_i r_i r_i^T=M_A+M_B.
\]
If $h=\lambda_{\min}(M)$, Lemma~\ref{lem:hub-gaussian} bounds the
conditional density by $4(2\pi)^{-3/4}h^{-3/4}$, uniformly in all
conditional means and in $P$.  It remains to average this bound.

For side $A$, set $R_A=(S_A^{-1/2}r_i)_{i\in A}$, so that
$R_AR_A^T=I_3$.  For any positive diagonal matrix $D$, compression
satisfies
\begin{equation}\label{eq:hub-compression}
 (R_ADR_A^T)^{-1/2}\le R_AD^{-1/2}R_A^T.
\end{equation}
Indeed the corresponding resolvent inequality follows from a Schur
complement, and integrating it against
$\pi^{-1}s^{-1/2}ds$ proves \eqref{eq:hub-compression}, using
$x^{-1/2}=\pi^{-1}\int_0^\infty s^{-1/2}(x+s)^{-1}\,ds$.
Since
$M_A=S_A^{1/2}[R_A\operatorname{diag}(T_i)R_A^T]S_A^{1/2}$,
taking traces and using \eqref{eq:hub-mixture-peak} gives
\[
 \mathbb E\lambda_{\min}(M_A)^{-1/2}
 \le\eta_3^{-1/2}\sum_{i\in A}
            |S_A^{-1/2}r_i|^2\mathbb ET_i^{-1/2}
 \le3\sqrt{2\pi}\eta_3^{-1/2}.
\]
The same bound holds for $M_B$.
The positive random variables
$a_0=\lambda_{\min}(M_A)$ and $b_0=\lambda_{\min}(M_B)$ are
independent, and $h\ge a_0+b_0\ge2\sqrt{a_0b_0}$.
Therefore, by independence and the monotonicity of positive moments,
\[
 \mathbb Eh^{-3/4}
 \le2^{-3/4}(\mathbb Ea_0^{-1/2})^{3/4}
                 (\mathbb Eb_0^{-1/2})^{3/4}
 \le2^{-3/4}(3\sqrt{2\pi}\eta_3^{-1/2})^{3/2}.
\]
Combining the two inequalities yields
\eqref{eq:hub-mixture-bound}.

For the final assertion, set $Y=X_0+\varepsilon P$.
Then its conditional density is $f_0(y-\varepsilon P)\le1$, and
\[
 aX_0^2+2X_0(L+\beta)+\lambda V^2+2\varepsilon LP
 =2YL+\lambda V^2+a(Y-\varepsilon P)^2
                         +2\beta(Y-\varepsilon P),
\]
which is in the stated class with $B=0$.
\end{proof}

\begin{remark}[Scope of the mixture condition]
Condition \eqref{eq:hub-mixture-peak} bounds the average peak of the
Gaussian components and implies that each marginal density is at most
one.  For centered Gaussian scale mixtures, meaning that each $M_i$
is deterministic, it is equivalent to the cap-one assumption: the
maximum is attained at $M_i$ and equals
$(2\pi)^{-1/2}\mathbb ET_i^{-1/2}$.
Thus Theorem~\ref{thm:hub-gaussian-mixtures} covers every cap-one
centered Gaussian scale mixture, even when its component variances
have no positive lower bound.  A mixture of many narrow components
separated in location can have a bounded final density while violating
\eqref{eq:hub-mixture-peak}; such laws are not covered.
\end{remark}

\begin{corollary}[Laplace light coordinates]\label{cor:hub-laplace}
Let the light coordinates have independent translated Laplace
densities $e^{-2|x-m_i|}$.  For the adapted forms of
Theorem~\ref{thm:hub-gaussian-mixtures}, no spread condition is needed,
and
\[
 \|p_Q\|_\infty\le12\pi^{-3/4}\Gamma(1/4)<19.
\]
\end{corollary}

\begin{proof}
The Laplace law has the representation
$m_i+\sqrt{T_i}G_i$ with $T_i$ exponential of rate two:
its characteristic function is
$\mathbb E e^{-t^2T_i/2}=4/(4+t^2)$.
For the covariance $M=\sum_iT_i r_i r_i^T$, the same resolvent
compression argument, now with exponent $3/4$, gives
\[
 \mathbb E\lambda_{\min}(M)^{-3/4}
 \le\mathbb E\operatorname{tr}M^{-3/4}
 \le\sum_i|r_i|^2\mathbb ET_i^{-3/4}
 =3\,2^{3/4}\Gamma(1/4).
\]
Average \eqref{eq:hub-gaussian-bound}.  No partition of the frame is
needed because the full negative $3/4$ moment exists here.
\end{proof}

\subsection{A limitation of density-cap arguments for adapted hubs}

\begin{proposition}[Conditional information is not controlled by marginal caps]
\label{prop:hub-cap-obstruction}
For every sufficiently large integer $m$, there is a joint law
$(L,V,P)$ whose joint density and all its coordinate marginal
densities are bounded by one, and a conditional hub density
$g(y\mid P)\le1$ of unit mass, such that
\[
 \|p_{2YL+V^2}\|_\infty\ge\frac14\log m-\frac12\log9.
\]
Here coordinate marginals include all one- and two-coordinate
projections of this three-dimensional vector.  The constructed joint
law is not asserted to be an orthogonal projection of a product law.
\end{proposition}

\begin{proof}
Let $\delta=1/m$.  Partition $[0,1]\times[1,2]$ into squares
\[
 C_{jk}=[j\delta,(j+1)\delta]\times
                    [1+k\delta,1+(k+1)\delta],\qquad0\le j,k<m.
\]
Assign to each square a distinct disjoint unit interval $I_{jk}$
on the $P$ axis.  Put joint density one on
$\bigcup_{j,k}C_{jk}\times I_{jk}$ and zero elsewhere.
Its total integral is $m^2\delta^2=1$.
The $(L,V)$ marginal is uniform on the rectangle; the $(L,P)$ and
$(V,P)$ densities are $\delta$ on their supports; and the $P$ density
is $\delta^2$.  The one-dimensional $L,V$ densities are one.

Write $(\ell_j,v_k)$ for the center of $C_{jk}$.
When $P\in I_{jk}$, let $Y$ be uniform on the unit interval centered
at $-v_k^2/(2\ell_j)$.  This kernel depends only on $P$.
For cells whose lower $L$ endpoint is at least $8\sqrt\delta$,
the root $-v^2/(2\ell)$ stays in this unit interval throughout the
cell.  The mean-value theorem bounds its distance from the center
root by
\[
 \frac\delta{\ell_{\min}}+
             \frac\delta{\ell_{\min}^2}
 \le\frac{\sqrt\delta}{8}+\frac1{64}<\frac12.
\]
The same remains true for the root $(s-v^2)/(2\ell)$ when
$|s|\le\sqrt\delta$, since the added displacement is at most
$1/16$.
Let $b_m$ be the first grid boundary at least $8\sqrt\delta$;
then $b_m\le9\sqrt\delta$ for large $m$.
For every such level $s$, the conditional density formula in $Y$
therefore gives
\[
 p_{2YL+V^2}(s)\ge
 \int_{b_m}^1\int_1^2\frac{dv\,d\ell}{2\ell}
 =\frac12\log(1/b_m)
 \ge\frac14\log m-\frac12\log9.
\]
This holds on a nonempty interval of levels, and hence proves an
essential-supremum lower bound.
\end{proof}

This example shows why the adapted problem requires more than caps
on the displayed joint and marginal densities: $P$ can encode a fine
$(L,V)$ cell and allow the hub to choose a separate center for that
cell.  It is an obstruction to a reduced proof strategy, not a
counterexample to the bounded-density conjecture for independent
coordinates.  The mixture theorem above uses both the conditional
Gaussian structure and the independence of two frame blocks, and
therefore retains information absent from this example.

\subsection{Collisions on shells and a logarithmic support dependence}

\begin{lemma}[A shell collision estimate]\label{lem:hub-shell}
Let $w,v,b$ be orthonormal, put $L=w\cdot X'$, $V=v\cdot X'$,
$B=b\cdot X'$, and let
$Y=\lambda(V-v_0)^2+\gamma B$, where $\lambda\ne0$.
For almost every $\ell$, let $h_\ell$ be the unnormalized density
of $Y$ on the slice $L=\ell$, restricted to
$r\le |V-v_0|\le K$, where $0<r<K$.  Then
\[
 \int_{\mathbb R}h_\ell(y)^2\,dy
 \le \frac8{|\lambda|}\log(K/r).
\]
No bound on the coordinate leverages is required.
\end{lemma}

\begin{proof}
Take an independent copy $(L',V',B')$ and put
$S=V+V'-2v_0$, $D=V-V'$, $E=B-B'$.
The square integral is the collision density at $(\ell,\ell,0)$
of $(L,L',Y-Y')$ for the two restricted measures, and
$Y-Y'=\lambda SD+\gamma E$.
First restrict to pairs on the same side of the fold $v_0$.
Then $2r\le |S|\le2K$.  For fixed $S=s$ put
$n=(\lambda s,\gamma)/\sqrt{\lambda^2s^2+\gamma^2}$.
On the product space of the two copies, the coefficient vectors of
\[
 L,\quad L',\quad (S+2v_0)/\sqrt2,\quad (n_1D+n_2E)/\sqrt2
\]
are orthonormal.  Their joint density is at most four by
\eqref{eq:hub-projection}.  Undoing the two factors $\sqrt2$ bounds
the density in the unscaled coordinates by two.  Restrictions of
the underlying measure can only decrease it.  Thus the contribution
of the same-side pairs is at most
\[
 \int_{2r\le |s|\le2K}
       \frac{2\,ds}{\sqrt{\lambda^2s^2+\gamma^2}}
 \le\frac4{|\lambda|}\log(K/r).
\]
For opposite-side pairs fix $D=d$ instead.  Now
$2r\le |D|\le2K$, and the four orthogonal directions are
$L,L',D/\sqrt2,(n_1S+n_2E)/\sqrt2$, with
$n=(\lambda d,\gamma)/\sqrt{\lambda^2d^2+\gamma^2}$.
The same calculation gives the remaining four units.
These collision identities can first be applied to mollified
densities and bounded restrictions; disintegration and monotone
convergence then give the asserted almost-everywhere statement.
\end{proof}

\begin{theorem}[Logarithmic dependence on bounded light supports]
\label{thm:hub-shell-support}
Suppose each light coordinate has density at most one and is
supported in an interval of length at most $K$.  In the normal form
at the start of this section suppose $a=0$ and the light diagonal
coefficients vanish.  Then, uniformly in all cross terms, affine
terms, and the hub law,
\[
 \|p_Q\|_\infty
 \le \frac12+\frac{C}{\sqrt{|\lambda|}}\,
              \sqrt{1+\log(1+K^2|\lambda|)}
\]
for an absolute constant $C$.
\end{theorem}

\begin{proof}
We have $K\ge1$.  On every frozen slice the light polynomial,
after completing the square, has the form in
Lemma~\ref{lem:hub-shell}, up to a constant; its direction $b$,
coefficient $\gamma$, and fold location may depend on the frozen
value.  A missing $b$ direction can be supplied by an extra,
independent coordinate whose coefficient in $Y$ is zero.
The exponential-tilt proof of
Proposition~\ref{prop:hub-light-intervals}, applied just in direction
$v$, gives
\[
 \rho_{L,V}(\ell,u)
 \le 2\exp\{-2(|u-\mathbb EV|-2K)_+^2/K^2\}.
\]
Consequently the slice mass on
$|V-\mathbb EV|>2K+t$ is at most
$\sqrt{2\pi}K e^{-2t^2/K^2}$.

Put $W=2K+t$ and $D_0=|v_0-\mathbb EV|$.
On the remaining part of the slice and outside the fold slab of
radius $r\le W$, the values of $|V-v_0|$ lie between
$\max(r,D_0-W)$ and $D_0+W$.  The ratio of these endpoints is at
most $3W/r$: if $D_0\le2W$ use the upper bound $3W$, and otherwise
use the ratio bound three.  The fold slab has slice mass at most
$4r$, since $\rho_{L,V}\le2$.  Cauchy--Schwarz and
Lemma~\ref{lem:hub-shell} therefore give, for every Borel set $S$,
\[
 \mu_{\ell_0,\ell}(S)
 \le4r+\sqrt{\frac8{|\lambda|}\log(3W/r)}\,\sqrt{|S|}
                  +\sqrt{2\pi}K e^{-2t^2/K^2}.
\]
Here and below harmless changes in the polynomial's additive
constant do not affect the estimate.

Write $x=|S|$ and $\mu=|\lambda|$.  If $x<\mu$, take
\[
 r=\sqrt{x/\mu},\qquad
 t=K\sqrt{\tfrac12\log_+
                       (\sqrt{2\pi}K\sqrt{\mu/x})}.
\]
Then $r<1\le K\le W$, and the last term is at most
$\sqrt{x/\mu}$.  For an absolute $C_0$ this proves
\[
 \mu_{\ell_0,\ell}(S)
 \le\frac{C_0}{\sqrt\mu}
       \sqrt{x\log\!\left(e+\frac{1+K^2\mu}{x}\right)}.
\]
Indeed $\log(3W/r)\le
C[1+\log(K\sqrt{\mu/x})]$, since the extra square root of a logarithm
in $W$ contributes at most the same expression.  If $x\ge\mu$,
the displayed bound follows instead from the total slice mass
$\rho_L(\ell)\le\sqrt2$, after enlarging $C_0$.

This last displayed expression defines an increasing concave
profile for all $x>0$, with value zero at zero.
To check concavity, $x\mapsto x\log(e+A/x)$ is increasing and
concave for every $A>0$, as direct differentiation shows;
composition with the increasing concave square root preserves both
properties.  Finally, substituting $x=2e^{-u}$ shows
\[
 \int_0^2\sqrt{x\log(e+A/x)}\,\frac{dx}{x}
 \le C_1\sqrt{1+\log(1+A)}.
\]
For example use
$\log(e+(A/2)e^u)\le u+\log(e+A/2)$ and integrate
$e^{-u/2}\sqrt{u+\log(e+A/2)}$.
Lemma~\ref{lem:hub-profile} proves the theorem.
\end{proof}

\begin{proposition}[A weighted collision criterion]
\label{prop:hub-weighted-collision}
In the notation of Lemma~\ref{lem:hub-shell}, set
\[
 G=\sqrt{\lambda^2
       (|V-v_0|+|V'-v_0|)^2+\gamma^2},
 \qquad Z=(Y-Y')/G.
\]
Suppose that, uniformly in the frozen parameters, the
disintegrated collision density of $(L,L',Z)$ at
$(\ell,\ell,0)$ is at most $C_0$ for almost every $\ell$.
More explicitly, in each of the two branch classes used in the proof
of Lemma~\ref{lem:hub-shell}, assume that the nonnegative density
$F_\ell(s)$ in the remaining shell coordinate at $Z=0$
satisfies $\int F_\ell(s)\,ds\le C_0$.
Then the density restricted to $|V-v_0|\ge r$ satisfies
\[
 \int h_\ell(y)^2\,dy
 \le\frac8{|\lambda|}
              \log\!\left(1+\frac{C_0}{8r}\right).
\]
In particular, if this hypothesis holds uniformly and
$|\lambda|\ge1$, the corresponding zero-diagonal one-hub forms
have uniformly bounded densities, with a constant depending only
on $C_0$.
\end{proposition}

\begin{proof}
In the same-side class use $s=S$ and in the opposite-side class
use $s=D$.  The orthogonal projection calculation above gives
$0\le F_\ell(s)\le2$, even though its direction depends on $s$.
The collision integral in either class is at most
$|\lambda|^{-1}\int_{|s|\ge2r}F_\ell(s)|s|^{-1}\,ds$.
The maximum under the constraints $F\le2$ and $\int F\le C_0$
is obtained by filling the intervals closest to the two endpoints
$\pm2r$.  Its value is
$4|\lambda|^{-1}\log(1+C_0/(8r))$.
Adding the two classes proves the first assertion.
The fold slab again has mass at most $4r$.  Hence
\[
 \mu_{\ell_0,\ell}(S)
 \le4r+\sqrt{\frac{8|S|}{|\lambda|}
                           \log(1+C_0/(8r))}.
\]
For $|\lambda|\ge1$, choose $r=\sqrt{|S|}$ near zero and use the
trivial slice mass bound away from zero.  An increasing concave
majorant is $C(C_0)\sqrt{x\log(e+(1+C_0)^2/x)}$.
Its integral against $dx/x$ on $(0,2)$ is finite, so
Lemma~\ref{lem:hub-profile} applies.
\end{proof}

The weighted collision hypothesis is an additional nonlinear
condition.  The flat projection density bound proves the
pointwise bound $F_\ell\le2$ but does not prove its integral
bound.  Thus the proposition gives a sufficient condition whose
validity for arbitrary product densities remains open.

\subsection{Averaging the hub quadratic phase}

\begin{lemma}[A parabola extension estimate]\label{lem:hub-strichartz}
If $0\le f\le1$ and $\int f=1$, and
$g(s,\tau)=\int f(x)e^{isx^2+i\tau x}\,dx$, then
\[
 \iint_{\mathbb R^2}|g(s,\tau)|^6\,ds\,d\tau
 \le \frac{4\pi^3}{\sqrt3}.
\]
The same bound holds if $f$ in the definition of $g$ is replaced
by a complex function dominated in absolute value by $f$.
\end{lemma}

\begin{proof}
The cube $g^3$ is the Fourier transform of the pushforward of
$f(x_1)f(x_2)f(x_3)\,dx$ under
$T(x)=(x_1+x_2+x_3,x_1^2+x_2^2+x_3^2)$.
If $r$ is the distance from $x$ to the line spanned by
$(1,1,1)$, the coarea Jacobian is $2\sqrt3\,r$.
Each nonempty regular fiber is a circle of radius $r$.
Its image density $m$ is consequently at most
$\pi/\sqrt3$, since the product of the three densities is at
most one.  It has integral one.  Plancherel gives
\[
 \iint |g|^6=(2\pi)^2\int m^2
 \le(2\pi)^2(\pi/\sqrt3)\int m.
\]
For a complex amplitude, the absolute value of its pushforward
is bounded by the same nonnegative density $m$, so the argument
is unchanged.  This coarea proof is the elementary
one-dimensional parabola case of the Fourier extension estimates
associated with \citet{Strichartz1977,Foschi2007,HZ2006}.
\end{proof}

\begin{theorem}[A degenerate hub with a quadratic partner]
\label{thm:hub-degenerate}
Let $l,u$ be linearly independent, let
$L=l\cdot X'$, $U=u\cdot X'$, and put
$G=\operatorname{Gram}(l,u)$.
Assume $(l_i,u_i)G^{-1}(l_i,u_i)^T\le1/2$ for every $i$.
For independent coordinates with density at most one, consider
\[
 Q=aH^2+2H(L+\beta)+\sigma L^2+\lambda U^2
                                      +\sum_i\phi_i(X_i),
\]
where the $\phi_i$ are arbitrary quadratic polynomials.
If $\sigma\lambda\ne0$ and $a\sigma\ne1$, define
\[
 \begin{split}
 C_6&=(4\pi^3/\sqrt3)^{1/6},\\
 B_0&=(4\pi)^{5/6}C_6
 |\lambda|^{-1/2}|\sigma|^{-1/6}|1-a\sigma|^{-1/6}
               (\det G)^{-5/12}|u|^{-1/6}.
 \end{split}
\]
Then, for every $T>0$,
\begin{equation}\label{eq:hub-strichartz-block}
 \int_T^{2T}|\mathbb E e^{itQ}|\,dt\le B_0T^{-1/6},
\end{equation}
and
\[
 \|p_Q\|_\infty
 \le\frac7\pi
 \left(\frac{B_0}{6(1-2^{-1/6})}\right)^{6/7}.
\]
\end{theorem}

\begin{proof}
The leverage assumption also implies
$u_i^2/|u|^2\le1/2$.  Absorb all linear hub terms into $\beta$.
Write $\phi_i(x)=d_ix^2+b_ix+c_i$ and
$\psi_i^{(t)}(\theta)=\int f_i(x)e^{itd_ix^2+i\theta x}\,dx$.
For a fixed hub value $x$ use the exact identity
\[
 ax^2+2xL+\sigma L^2
   =\sigma(L+x/\sigma)^2+(a-1/\sigma)x^2.
\]
Two Fresnel linearizations, with Gaussian damping followed by
dominated convergence, therefore give
\[
 |\varphi_Q(t)|\le
 \frac1{4\pi|t|\sqrt{|\sigma\lambda|}}
 \int_{\mathbb R}
 |g_H(t(a-1/\sigma),2t\beta+\zeta/\sigma)|K_t(\zeta)\,d\zeta,
\]
where
$K_t(\zeta)=\int\prod_i
 |\psi_i^{(t)}(\zeta l_i+\xi u_i+tb_i)|\,d\xi$.
Lemma~\ref{lem:hub-frame-integral}, first in the $u$ direction
and then in both directions, yields
\[
 \|K_t\|_\infty\le4\pi/|u|,
 \quad \|K_t\|_1\le(4\pi)^2(\det G)^{-1/2},
 \quad
 \|K_t\|_{6/5}\le
 (4\pi)^{11/6}(\det G)^{-5/12}|u|^{-1/6}.
\]
Apply H\"older with exponents $6$ and $6/5$ on
$[T,2T]\times\mathbb R$.  The map
\[
 (t,\zeta)\longmapsto
 (t(a-1/\sigma),\,2t\beta+\zeta/\sigma)
\]
has Jacobian $|1-a\sigma|/\sigma^2$.
Lemma~\ref{lem:hub-strichartz} bounds the first factor by
$C_6(\sigma^2/|1-a\sigma|)^{1/6}$.
The second factor is at most
$T^{5/6}(4\pi)^{11/6}(\det G)^{-5/12}|u|^{-1/6}$.
Multiplying these factors and $1/(4\pi T\sqrt{|\sigma\lambda|})$
proves \eqref{eq:hub-strichartz-block}.
Since $Q$ is real, the negative-frequency bound is identical.
The full Fourier integral is at most
\[
 2T_0+\frac{2B_0T_0^{-1/6}}{1-2^{-1/6}}.
\]
Its minimum occurs at
$T_0=[B_0/(6(1-2^{-1/6}))]^{6/7}$.
Fourier inversion proves the density estimate.
\end{proof}

\begin{theorem}[Uniformity through zero partner curvature]
\label{thm:hub-degenerate-uniform}
In Theorem~\ref{thm:hub-degenerate}, assume $a=0$.
There is a finite bound depending only on
$|\lambda|$, $\det G$, and $|u|$, uniform in every
$\sigma\in\mathbb R$, in the affine terms, and in the light
diagonal terms.  One explicit bound is
\[
 \|p_Q\|_\infty\le
 1+\frac{13}{4}C_H+
 C_{\mathrm{av}}|\lambda|^{-1/2}
                         (\det G)^{-5/12}|u|^{-1/6},
 \qquad
 C_H=3\,2^{4/3}(|\lambda|\det G)^{-1/3},
\]
where
\[
 C_{\mathrm{av}}=
 \frac{(4\pi)^{5/6}C_6}{\pi}
 \left\{\frac2{1-2^{-1/6}}
             +\frac{2}{3(2^{11/6}-1)}\right\}.
\]
\end{theorem}

\begin{proof}
When $\sigma=0$, Proposition~\ref{prop:hub-localized} with empty
localization set, after changing the two frame coordinates,
gives the bound $C_H$.  Suppose $\sigma\ne0$.
Set $Q_1=Q-\sigma L^2$ and
$Q_2=Q-\sigma(L+\beta)^2$.
Each has the preceding zero-curvature form and density at most
$C_H$, since their difference consists only of affine light terms.
Fix $\delta\ge|\sigma|$ and put
$\Lambda=\sqrt{\delta/|\sigma|}\ge1$.
On $|L|\le\Lambda$, comparison with $Q_1$ bounds a window of
radius $\delta$ by $4C_H\delta$.
On $|L|>\Lambda$, $|L+\beta|\ge\Lambda/2$, conditioning on
$X'$ and using the affine hub slope gives probability at most
$\delta/|L+\beta|\le2\delta$.
On the remaining event,
$|Q-Q_2|\le|\sigma|(L+\beta)^2<\delta/4$,
so comparison with $Q_2$ gives at most $(5/2)C_H\delta$.
Consequently
\[
 \sup_s\frac{\mathbb P(|Q-s|\le\delta)}{2\delta}
 \le1+\frac{13}{4}C_H.
\]

To pass to smaller scales, put $h=|\sigma|$.
The preceding theorem gives an integrable characteristic function.
For the uniform probability density $U_h$ on $[-h,h]$,
Fourier inversion and $0\le1-\sin z/z\le\min(2,z^2/6)$ imply
\[
 \|p_Q-p_Q*U_h\|_\infty
 \le\frac1{2\pi}\int |\varphi_Q(t)|
                       (1-\operatorname{sinc}(ht))\,dt.
\]
Use \eqref{eq:hub-strichartz-block} on blocks
$T=2^j/h$, $j\in\mathbb Z$.
On such a block the last factor is at most
$\min(2,(2/3)4^j)$.  The geometric sums over $j\ge0$ and
$j\le-1$ are respectively bounded by
$2/(1-2^{-1/6})$ and $2/[3(2^{11/6}-1)]$.
Thus the error is at most
$B_0h^{1/6}$ times their sum divided by $\pi$.
When $a=0$, the factor $|\sigma|^{-1/6}$ in $B_0$
cancels $h^{1/6}$ exactly.  The estimate for $p_Q*U_h$
already follows from the window bound at $\delta=h$.
Adding the two bounds proves the result.
\end{proof}

\begin{corollary}[A normalized degenerate hub]
\label{cor:hub-degenerate-normalized}
If $w,v$ are orthonormal and $w_i^2+v_i^2\le1/2$, then
\[
 4H(w\cdot X')+2(w\cdot X')^2+2(v\cdot X')^2
                         +\sum_i\phi_i(X_i)+b_HH
\]
has density less than $18$, for every choice of the independent
coordinate laws with density at most one.
\end{corollary}

\begin{proof}
Apply Theorem~\ref{thm:hub-degenerate} with $l=2w$, $u=v$,
$a=0$, $\sigma=1/2$, and $\lambda=2$.
Then $\det G=4$, $|u|=1$, and the exact expression for $B_0$
is less than $7.5$.  The resulting exact density bound is
less than $18$.
\end{proof}

For a quadratic polynomial confined to $\operatorname{span}(w,v)$,
completing the square in $V$ produces
$\lambda'(V+\kappa L+\gamma)^2+\sigma L^2$ and affine terms.
The replacement $u=v+\kappa w$ preserves
$\det\operatorname{Gram}(w,u)=1$ and the relative leverage
$w_i^2+v_i^2$.  Hence the uniform theorem applies whenever the
transverse coefficient $|\lambda'|$ stays bounded below.
For $a\ne0$, the bound retains its factors
$|\sigma|^{-1/6}|1-a\sigma|^{-1/6}$ and diverges as the
corresponding curvatures vanish.

\subsection{Hub diagonals and square-root profiles with random centers}

Here $H$ is independent of $(t,R)$, has density $f_H\le1$, and
\[
 Q=aH^2+2Ht+R.
\]
The random variable $t$ may be a shifted linear projection of the
light coordinates.  We use the following explicitly joint
hypothesis: there are a random center $C$ coupled with $t$ and
a constant $K$ such that for every nonnegative Borel function
$\Phi$,
\begin{equation}\label{eq:hub-sq}
 \mathbb E\Phi(t,R)\le
 K\mathbb E\int_{\mathbb R}
                         \Phi(t,r)|r-C|^{-1/2}\,dr.
\end{equation}
For example, if
$R=\lambda(V-D_0)^2+C$ and $V$ has conditional density at most
$\kappa$ given $(t,D_0,C)$, then \eqref{eq:hub-sq} holds with
$K=\kappa/\sqrt{|\lambda|}$, by the two-root change of variables.
The auxiliary variables used to represent the right-hand side of
\eqref{eq:hub-sq} are sampled independently of $H$.

\begin{lemma}[A scalar discriminant bound]\label{lem:hub-sq-scalar}
For a probability density $0\le f\le1$ and
$\Delta=t^2+a s\ne0$,
\[
 \int_{\mathbb R}f(x)|s-ax^2-2tx|^{-1/2}\,dx
 \le2^{7/4}|\Delta|^{-1/4}.
\]
If $a=0$ and $t\ne0$, the sharper bound is $2|t|^{-1/2}$.
\end{lemma}

\begin{proof}
For $a\ne0$, put $A=|a|$, $d=|\Delta|$, and
$z=A(x+t/a)/\sqrt d$.
The sublevel condition $|s-ax^2-2tx|<u$ becomes
$|z^2-\operatorname{sgn}\Delta|<Au/d$.
If $\Delta>0$, the length of this set in the $z$ variable is
\[
 2\{\sqrt{1+r}-\sqrt{(1-r)_+}\},\qquad r=Au/d,
\]
which is at most $2\sqrt2\,r$.  For $r\le1$ rationalize
the difference and use
$\sqrt{1+r}+\sqrt{1-r}\ge\sqrt2$; for $r\ge1$ use
$\sqrt{1+r}\le\sqrt2\,r$.
If $\Delta<0$, its length is $2\sqrt{(r-1)_+}\le r$,
so the same bound holds.  After returning to $x$,
\[
 |\{x:|s-ax^2-2tx|<u\}|
                      \le 2\sqrt2\,u/\sqrt d.
\]
Layer cake and $f\le1$, $\int f=1$ yield
\[
 \int f|s-ax^2-2tx|^{-1/2}
 \le\int_0^\infty
       \min(1,2\sqrt2\,v^{-2}/\sqrt d)\,dv
 =2^{7/4}d^{-1/4}.
\]
For $a=0$ the sublevel length is exactly $u/|t|$,
and the same integral gives $2|t|^{-1/2}$.
\end{proof}

\begin{theorem}[Composition with a random center]
\label{thm:hub-sq-composition}
Under \eqref{eq:hub-sq}, for almost every $s$,
\begin{equation}\label{eq:hub-sq-composition}
 p_Q(s)\le1+2^{7/4}K\,
                    \mathbb E|t^2+a(s-C)|^{-1/4}.
\end{equation}
If $C=\sigma_c t^2+\eta_1t+\eta_0+C'$ for deterministic
$\sigma_c,\eta_1,\eta_0$, and the conditional density of $t$
given $C'$ is bounded by $D$, then, provided
$D_H=1-a\sigma_c\ne0$,
\[
 \|p_Q\|_\infty
 \le1+2^{13/4}K\sqrt D\,|D_H|^{-1/4}.
\]
At $a=0$ the first bound requires no information about the
random center beyond \eqref{eq:hub-sq}.
\end{theorem}

\begin{proof}
For fixed $(t,r)$ the roots of $ax^2+2tx+r=s$ have derivative
$2|ax+t|$.  Roots with $|ax+t|\ge1$ contribute at most one
in total.  Let
\[
 \mathcal K_s(t,r)=
 \sum_{\substack{x:\,ax^2+2tx+r=s\\|ax+t|<1}}
                         \frac{f_H(x)}{2|ax+t|}.
\]
Changing $r$ at fixed $t$ on each monotonicity branch gives
\[
 \int \mathcal K_s(t,r)\Psi(r)\,dr
 =\int f_H(x)\mathbf1_{\{|ax+t|<1\}}
                    \Psi(s-ax^2-2tx)\,dx.
\]
Apply \eqref{eq:hub-sq} to $\mathcal K_s$, use this identity,
and drop the indicator.  Lemma~\ref{lem:hub-sq-scalar},
with its level parameter replaced by $s-C$, proves
\eqref{eq:hub-sq-composition}.
The formulas also follow by smoothing in $s$, which avoids
exceptional critical values; Tonelli then justifies all the
nonnegative integrals.

Conditionally on $C'$, the discriminant in the second assertion
is a quadratic polynomial in $t$ with leading coefficient $D_H$.
For real roots, factor it and use Cauchy--Schwarz.
For conjugate nonreal roots, its absolute value is at least
$|D_H|(t-r)^2$, where $r$ is their common real part.
In either case
\[
 \mathbb E\bigl(|t^2+a(s-C)|^{-1/4}\mid C'\bigr)
 \le |D_H|^{-1/4}
             \sup_r\mathbb E(|t-r|^{-1/2}\mid C')
 \le2\sqrt2\,\sqrt D\,|D_H|^{-1/4}.
\]
The final bound is the bathtub integral over the interval of
length $1/D$ centered at $r$.  Multiplication of the constants
proves the result.
\end{proof}

\begin{proposition}[Randomness at the degenerate center]
\label{prop:hub-sq-log}
Assume \eqref{eq:hub-sq}, the center representation, and the
conditional density cap $D$ in
Theorem~\ref{thm:hub-sq-composition}, now allowing $D_H=0$.
Suppose $\sigma_c\ne0$.
For a random variable $Z$ write
$\alpha_1(Z)=\sup_{s\in\mathbb R,\,u>0}\mathbb P(|Z-s|\le u)/\sqrt u$.
With the convention that the second argument below is infinite
when $D_H=0$, put
\[
 A_*=\min\left\{\alpha_1(C'),
                    2\sqrt2\sqrt{|\sigma_c|/|D_H|}\right\}.
\]
If $A_*<\infty$, then
\[
 \|p_Q\|_\infty
 \le1+KD|\sigma_c|^{-1/2}
       \{12+2\log_+|\sigma_c|+4\log_+ A_*\}.
\]
In particular, center randomness with a finite
$\alpha_1(C')$ gives a finite bound at $D_H=0$.
\end{proposition}

\begin{proof}
Retain the indicator in the proof of
Theorem~\ref{thm:hub-sq-composition}, and condition on $(H,C')$.
The density cap $D$ reduces the remaining integral to
\[
 D\int_{|ax+\tau|<1}
 |s-c'-\eta_0-ax^2-(2x+\eta_1)\tau-\sigma_c\tau^2|^{-1/2}
                                                        \,d\tau.
\]
The interval has length two.  A quadratic with leading
coefficient $-\sigma_c$ and critical value $h_c$ has integral
over any such interval at most
\begin{equation}\label{eq:hub-quadratic-log-kernel}
 |\sigma_c|^{-1/2}
       \{8+2\log_+(|\sigma_c|/|h_c|)\}.
\end{equation}
To verify this bound, after translating and dividing by
$|\sigma_c|$ write the integrand as $|z^2-c|^{-1/2}$.
On $|z|\le2\sqrt{|c|}$ its integral over the whole region
is at most $\pi+2\operatorname{arcosh}2<6$.
Outside that region it is at most $2/(\sqrt3|z|)$.
On a set of length at most two the integral of this last
majorant is at most
$(4/\sqrt3)\log(1+1/(2\sqrt{|c|}))$.
Since this is at most
$(4/\sqrt3)\log(3/2)+(2/\sqrt3)\log_+(1/|c|)$,
\eqref{eq:hub-quadratic-log-kernel} follows.  The case $c=0$
is interpreted with an infinite right side.

The critical value here is $s-C'-q_c(H)$, where
\[
 q_c(x)=\eta_0+ax^2-(x+\eta_1/2)^2/\sigma_c,
\]
whose quadratic coefficient is $-D_H/\sigma_c$.
The variables $C'$ and $H$ are independent.
Convolution does not increase the concentration function, so
\[
 \mathbb P(|C'+q_c(H)-s|\le u)
 \le\min(1,A_*\sqrt u).
\]
For the bound from $q_c(H)$ use the elementary quadratic
sublevel length $2\sqrt{2u/|D_H/\sigma_c|}$ and $f_H\le1$.
Layer cake now gives
\[
 \mathbb E\log_+\frac1{|s-C'-q_c(H)|}
 \le\int_0^\infty\min(1,A_*e^{-v/2})\,dv
 \le2+2\log_+ A_*.
\]
Insert this and
$\log_+(|\sigma_c|/|h_c|)
\le\log_+|\sigma_c|+\log_+(1/|h_c|)$
into \eqref{eq:hub-quadratic-log-kernel}.
The near-root contribution is the asserted expression after
multiplication by $KD$, and the far roots cost at most one.
\end{proof}

\begin{corollary}[An independent transverse square]
\label{cor:hub-sq-independent}
Suppose $V=v\cdot X_S$ and $t=w\cdot X_T+\beta$ use disjoint
coordinate sets, and $D_0,C$ are arbitrary functions of $X_T$.
For
$Q=aH^2+2Ht+\lambda(V-D_0)^2+C$,
\[
 \|p_Q\|_\infty
 \le\frac{2^{13/4}}{|v|\sqrt{|\lambda|}}\,
        \alpha_1(t^2-aC)^{1/2}.
\]
When $a=0$ one has the unconditional bound
\[
 \|p_Q\|_\infty
 \le 2^{13/4}|\lambda|^{-1/2}|v|^{-1}|w|^{-1/2}.
\]
\end{corollary}

\begin{proof}
Condition first on $(H,X_T)$ and integrate the independent square.
Since $\|p_V\|_\infty\le\sqrt2/|v|$, the resulting density
is bounded by
\[
 \frac{\sqrt2}{|v|\sqrt{|\lambda|}}\,
 \mathbb E|s-aH^2-2Ht-C|^{-1/2}.
\]
There is no far-root term in this argument.
The scalar lemma bounds the expectation by
$2^{7/4}\mathbb E|t^2+a(s-C)|^{-1/4}$.
For any $Z$ with $\alpha_1(Z)<\infty$, layer cake gives
$\mathbb E|Z-s|^{-1/4}\le2\alpha_1(Z)^{1/2}$.
This proves the first displayed bound.
For $a=0$, use instead the scalar estimate
$2|t|^{-1/2}$ and the linear density cap
$\sqrt2/|w|$.  Its bathtub integral is
$2\sqrt2(\sqrt2/|w|)^{1/2}$, which yields the second bound.
\end{proof}

The square-root hypothesis is a statement about the joint law
with the hub partner.  A concentration estimate for the marginal
law of $R$ alone cannot replace it.  For example, if
$R=2Yt$ with an additional independent variable $Y$, then at
$a=0$ the complete form is $2(H+Y)t$; for suitable interval
laws this product has a logarithmically unbounded density.
Likewise, at $C=t^2/a$ and opposite signs of $a$ and $\lambda$,
the form $(aH+t)^2/a+\lambda V^2$ has a hyperbolic
two-dimensional singularity.  These examples have only two
effective quadratic directions and explain the determinant
and center-randomness requirements without contradicting the
rank-three target.
